\chapter{Big backgrounds under bounded mass}
\label{chap:bounded-mass-equations}

Throughout this chapter $X$ is a connected compact complex manifold of dimension $n$ with bounded mass, $\theta$ is a smooth big real $(1,1)$-form on $X$, and the model reference $\varphi\in\PSH(X,\theta)$ satisfies \eqref{eq:li-model-condition}. Semipositivity of $\theta$ is not assumed. Bounded mass and positive lower reference volume give a finite positive interval for masses in the relative class. The convergence and equation arguments below use these mass controls.

\clearpage
\section{The general-measure equations}
\label{sec:big-general-equation-setting}

Retain the manifold $X$, dimension $n$, background $\theta$ and reference $\varphi$ just specified. Fix a Hermitian form $\omega$ and a smooth positive volume form $\mathrm dV$ on $X$. Bounded mass means $v_+(\omega)<\infty$, where $v_+(\omega)$ and $v_-(\omega)$ are the supremum and infimum of the Monge--Amp\`ere masses of bounded $\omega$-plurisubharmonic potentials. The lower and upper relative volumes and the class $\mathcal E_{\mathrm{vol}}(X,\theta,\varphi)$ are those of \cref{chap:mass-bound-volumes}. We write $u\simeq\varphi$ when their difference is bounded off their common poles.

\Needspace{22\baselineskip}
\begin{theorem}
\label{thm:li-general-measure}
Assume that $X$ has bounded mass. Let $\theta$ be big and let $\varphi$ satisfy \eqref{eq:li-model-condition}; in particular,
$\underline{\operatorname{vol}}_\theta(\varphi)>0$. For every nonzero positive non-pluripolar Radon measure $\mu$:
\begin{enumerate}
\item There are $u\in\mathcal E_{\mathrm{vol}}(X,\theta,\varphi)$ and a unique $c>0$ such that $\mathrm{MA}_\theta(u)=c\mu$.
\item For each $\lambda>0$ there is a unique $u\in\mathcal E_{\mathrm{vol}}(X,\theta,\varphi)$ such that $\mathrm{MA}_\theta(u)=\mathrm{e}^{\lambda u}\mu$.
\end{enumerate}
If $\mu=f\,\mathrm{d}V$ with $0\leq f\in L^p(\mathrm{d}V)$ for some $p>1$, the solutions in both parts may be taken with $u\simeq\varphi$.
\end{theorem}

Once the positive-volume model reference is fixed, the metric positivity is automatic under bounded mass. Indeed, choose $C>0$ with $\theta\leq C\omega$. Since $\varphi\preceq0$, \cref{thm:volume-monotonicity} gives
\[
 0<\underline{\operatorname{vol}}_\theta(\varphi)
 \leq\underline{\operatorname{vol}}_{C\omega}(0)
 =C^n v_-(\omega).
\]
The proofs in \cref{sec:big-relative-projection,sec:big-model-general-measures,sec:big-model-lp} reconstruct all the assertions. They construct the rooftop between different reference classes before projecting a minimal-reference solution, maximize the weighted measure of supersolutions, and obtain the $L^p$ ordinary type by a relative moment estimate. The class $\mathcal E_{\mathrm{vol}}$ agrees with $\mathcal E_{\mathrm{env}}$ by \cref{thm:relative-equivalence}.

\section{Convergence under a fixed measure}
\label{sec:big-relative-foundations}
\label{sec:big-relative-continuity}

As in the semipositive case, a fixed dominating measure gives a common lower envelope after extraction, and then convergence in capacity.

\Needspace{20\baselineskip}
\begin{theorem}\label{thm:big-relative-fixed-continuity}
Suppose $u_j\in\mathcal E_\varphi$, $u_j\to u\in\PSH(X,\theta)$ in $L^1(X)$, and $\mu_{u_j}\leq\nu$ for one finite non-pluripolar measure $\nu$. Then there is a subsequence with
\[
 h=P_\theta\left(\inf_j u_j\right)\in\mathcal E_\varphi.
\]
Moreover, for the original sequence,
\[
 u\in\mathcal E_\varphi,\qquad
 u_j\longrightarrow u\text{ in capacity},\qquad
 \mu_{u_j}\rightharpoonup\mu_u\leq\nu.
\]
\end{theorem}
\begin{proof}
The proofs of \cref{lem:dominated-relative-rooftop-extraction}, \cref{prop:relative-common-envelope-criterion} and \cref{thm:relative-fixed-measure-continuity} apply, with three replacements: the envelope identity replaces \eqref{eq:relative-normalized-model-bound}, the lower bound $\ell$ of \eqref{eq:big-relative-mass-interval} replaces $\kappa$, and zero domination replaces \cref{lem:relative-zero-domination}. Their other ingredients do not use semipositivity of the background: finite rooftops and stretching envelopes of members of $\mathcal E_\varphi$ stay in $\mathcal E_\varphi$ (\cref{prop:relative-finite-rooftops,lem:relative-stretch-closure}), contact support holds by \cref{thm:big-general-envelope-contact}, and \cref{lem:fixed-npp-integrals,lem:monotone-psh-capacity,lem:weighted-extended-lsc,lem:qpsh-hartogs,thm:general-background-damping} hold for every smooth background.

In particular, after a common translation with $u_j,u\leq\varphi\leq0$ and an extraction with $\int_X|\mathrm{e}^{u_j}-\mathrm{e}^u|\,\mu_{u_j}\leq2^{-j}$, the finite rooftops $p_{j,N}=P_\theta(u_j,\ldots,u_N)$ satisfy, as in \eqref{eq:finite-rooftop-common-domination},
\begin{equation}\label{eq:big-finite-roof-domination}
 \mu_{p_{j,N}}\leq\nu,\qquad
 \int_X|\mathrm{e}^{p_{j,N}}-\mathrm{e}^u|\,\mu_{p_{j,N}}\leq2^{1-j},
\end{equation}
and the stretched rooftops $Q_B(p_{j,N},\varphi)\in\mathcal E_\varphi$ have masses at least $\ell$. The exponential test \eqref{eq:relative-rooftop-noncollapse-test} then gives the common lower envelope $h\in\mathcal E_\varphi$ of the first assertion. In the proof of \cref{prop:relative-common-envelope-criterion}, zero domination applies to the limit $p\in\mathcal E_\varphi$ of the tail envelopes, since $p\leq u\preceq\varphi$ and $\mu_p(\{p<u\})=0$; it gives $p=u$.
\end{proof}


\Needspace{20\baselineskip}
\section{Relative moments}
\label{sec:big-relative-first-moment}

For $w\in\mathcal E_\varphi$ and $u\leq w$, the relative first moment $\int_X(w-u)\,\mu_u$ measures how far $u$ lies below $w$, weighted by the measure of $u$. A uniform bound on these moments prevents a sequence from collapsing and keeps its limits in $\mathcal E_\varphi$. The proof compares rooftop masses with the fixed measure of a potential of the same type as $w$ whose mass is close to the upper volume $M$. For a real number $t$, write $t_+=\max(t,0)$.

\begin{proposition}\label{prop:big-relative-first-moment}\index{Relative full-mass class!first-moment criterion}
Fix $w\in\mathcal E_\varphi$. Suppose $u_j\in\mathcal E_\varphi$, $u_j\leq w$, and
\[
 K:=\sup_j\int_X(w-u_j)\,\mu_{u_j}<+\infty.
\]
Then $\sup_Xu_j$ is bounded below. If $u_j\to u\in\PSH(X,\theta)$ in $L^1(X)$, then $u\in\mathcal E_\varphi$.

The same conclusions hold if $u_j\leq w$ is replaced by a common upper bound $u_j\leq\varphi+C$ and the integral is replaced by $\int_X(w-u_j)_+\,\mu_{u_j}$.
\end{proposition}
\begin{proof}
We first treat the case $u_j\leq w$. Since $\overline{\operatorname{vol}}_\theta(w)=M$, we may choose $w'\sim w$ with $\mu_{w'}(X)>M-\ell/2$. Translate $w'$ so that $w'\leq w$, and put $p_j=P_\theta(u_j,w')$. These rooftops belong to $\mathcal E_\varphi$. Contact support and ordered contact imply
\begin{equation}\label{eq:big-energy-roof-bound}
 \int_X(w'-p_j)\,\mu_{p_j}
 \leq\int_X(w-u_j)\,\mu_{u_j}\leq K.
\end{equation}
Indeed, the integrand vanishes on the contact part $\{p_j=w'\}$, and on the remaining contact part $p_j=u_j<w'\leq w$.

\Needspace{8\baselineskip}
We shall use the following estimate for maxima. If $a\in\mathcal E_\varphi$, $a\leq w'$, and $t>0$, then
\begin{equation}\label{eq:big-energy-max-bound}
 \mu_a(\{a>w'-t\})
 \leq \mu_{w'}(\{a>w'-t\})+M-\mu_{w'}(X).
\end{equation}
To see this, put $m=\max(a,w'-t)$. It has the ordinary singularity type of $w'$, so $\mu_m(X)\leq M$. On $\{a>w'-t\}$, locality gives $\mu_m=\mu_a$. On the complementary set, where $m=w'-t$, ordered contact gives $\mu_m\geq\mu_{w'}$. Integrating over the two sets proves \eqref{eq:big-energy-max-bound}.

Suppose that a subsequence had $\sup_Xu_j\to-\infty$; then also $\sup_Xp_j\to-\infty$. Splitting the mass of $p_j$ at the level $w'-s$, $s>0$, and using \eqref{eq:big-energy-roof-bound} and \eqref{eq:big-energy-max-bound}, we obtain
\[
 \ell\leq\mu_{p_j}(X)
 \leq\frac Ks+\mu_{w'}(\{p_j>w'-s\})+M-\mu_{w'}(X).
\]
For fixed $s$, the middle term tends to zero as $j\to\infty$, since $w'$ is finite off a pluripolar set and $\mu_{w'}$ is finite and non-pluripolar. Letting $j\to\infty$ and then $s\to\infty$ gives $\ell\leq M-\mu_{w'}(X)<\ell/2$, a contradiction. This proves the first assertion.

\Needspace{8\baselineskip}
Now suppose $u_j\to u$ in $L^1$. Then $u\leq w\preceq\varphi$. Fix $b>1$ and suppose that $Q_b(u,w)\equiv-\infty$. Since $w-w'$ is bounded, translating competitors by a constant shows that $Q_b(u,w')\equiv-\infty$ as well. Keep the same $p_j$ and set
\[
 h_j=Q_b(p_j,w')\in\mathcal E_\varphi.
\]
We claim that $\sup_Xh_j\to-\infty$. Otherwise, since $h_j\leq p_j\leq w'$, \cref{thm:normalized-compactness} would give a subsequence converging in $L^1$ to some $h\not\equiv-\infty$. Passing to the limit in
\[
 b^{-1}h_j+(1-b^{-1})w'\leq p_j\leq u_j,
\]
first almost everywhere and then everywhere by \cref{lem:qpsh-almost-everywhere-order}, would give a competitor for $Q_b(u,w')$, a contradiction.

On the contact set of $h_j$ one has
$p_j=b^{-1}h_j+(1-b^{-1})w'$. Thus $h_j\leq w'-bs$ implies $p_j\leq w'-s$ there. Contact support and \eqref{eq:big-energy-roof-bound} therefore imply
\[
 \mu_{h_j}(\{h_j\leq w'-bs\})\leq b^nK/s.
\]
Applying \eqref{eq:big-energy-max-bound} with $a=h_j$ and $t=bs$ on the complementary set, we obtain
\[
 \ell\leq b^nK/s+
     \mu_{w'}(\{h_j>w'-bs\})+M-\mu_{w'}(X).
\]
\Needspace{7\baselineskip}
Let first $j\to\infty$ and then $s\to\infty$. The choice of $w'$ again gives $\ell<\ell/2$, a contradiction. Therefore $Q_b(u,w)\not\equiv-\infty$ for every $b>1$. The resulting convex inequalities and \cref{thm:volume-monotonicity,prop:volume-addition} give
\[
 (1-b^{-1})^n M\leq\overline{\operatorname{vol}}_\theta(u)\leq M.
\]
Letting $b\to\infty$ shows that the upper volumes of $u$ and $\varphi$ are equal; hence $u\in\mathcal E_\varphi$.

\Needspace{6\baselineskip}
For the last assertion put $q_j=P_\theta(u_j,w)$. Contact gives
\[
 \int_X(w-q_j)\,\mu_{q_j}
 \leq\int_X(w-u_j)_+\,\mu_{u_j}\leq K.
\]
The case already treated bounds $\sup_Xq_j$ from below, and therefore also $\sup_Xu_j\geq\sup_Xq_j$. If $u_j\to u$ in $L^1$, then, since $q_j\leq w$, a subsequence of the $q_j$ converges in $L^1$ to some $q\not\equiv-\infty$, and the case already treated gives $q\in\mathcal E_\varphi$. Since $q\leq u\leq\varphi+C$, upward stability gives $u\in\mathcal E_\varphi$ as well.
\end{proof}

A common lower bound in $\mathcal E_\varphi$ and a weighted first moment of the difference from the limit tending to zero give convergence in capacity: rooftops of the tails identify the limit.

\Needspace{13\baselineskip}
\begin{lemma}\label{lem:big-weighted-common-bound}
Let $a_j\in\PSH(X,\theta)$ have a common upper bound and converge in $L^1(X)$ to $a\in\PSH(X,\theta)$ with $a\preceq\varphi$. Suppose there is $h\in\mathcal E_\varphi$ such that $h\leq a_j$ for every $j$. If, for some $\lambda>0$,
\[
 \int_X|a_j-a|\mathrm{e}^{\lambda a}\,\mu_{a_j}\longrightarrow0,
\]
then $a\in\mathcal E_\varphi$ and $a_j\to a$ in capacity.
\end{lemma}
\begin{proof}
By \cref{lem:qpsh-almost-everywhere-order}, the lower bound passes to the limit, so $h\leq a\preceq\varphi$, and upward stability gives $a\in\mathcal E_\varphi$. Start with any subsequence and extract a further subsequence along which the integrals in the hypothesis are at most $2^{-j}$ and $a_j\to a$ almost everywhere for smooth volume. For $N\geq j$ let $p_{j,N}=P_\theta(a_j,\ldots,a_N)$; these rooftops are not identically $-\infty$ because $h$ is a competitor. Partitioning the contact set according to the first obstacle touched gives
\[
 \int_X\min(1,|p_{j,N}-a|)\mathrm{e}^{\lambda a}\,\mu_{p_{j,N}}
 \leq\sum_{k=j}^N\int_X|a_k-a|\mathrm{e}^{\lambda a}\,\mu_{a_k}
 \leq2^{1-j}.
\]
The decreasing limits $p_j=P_\theta(\inf_{k\geq j}a_k)$ satisfy $h\leq p_j\leq a$, the second inequality by the almost-everywhere convergence, and hence so does their increasing regularized limit $p$. Upward stability puts the $p_j$ and $p$ in $\mathcal E_\varphi$, although the finite rooftops need not lie there.

Both $p_{j,N}\downarrow p_j$ and $p_j\uparrow p$ hold in capacity by \cref{lem:monotone-psh-capacity}, since $h$ is a common lower bound. Apply \cref{lem:weighted-extended-lsc} to this rooftop estimate, with the bounded weights $\min(1,|p_{j,N}-a|)\mathrm{e}^{\lambda a}$ and then $\min(1,|p_j-a|)\mathrm{e}^{\lambda a}$. These weights converge in capacity: discard the deep sublevel sets of $h$ and $a$, and use continuity on the remaining finite ranges. The result is
\[
 \int_X\min(1,|p-a|)\mathrm{e}^{\lambda a}\,\mu_p=0.
\]
Since the exponential is positive off a pluripolar set, which $\mu_p$ does not charge, $\mu_p(\{p<a\})=0$, and zero domination gives $p=a$. The upper envelopes $(\sup_{k\geq j}a_k)^*$ decrease to $a$. Squeezing the extracted $a_j$ between these envelopes and $p_j$ proves convergence in capacity. Since the initial subsequence was arbitrary, the conclusion holds for the original sequence.
\end{proof}

Without a given lower bound, the first-moment criterion produces one from the rooftops of the tails.

\Needspace{20\baselineskip}
\begin{theorem}\label{thm:big-relative-self-measure-convergence}
Let $u_j,u\in\mathcal E_\varphi$, with $u_j\to u$ in $L^1(X)$. If
\[
 \int_X|u_j-u|\,\mu_{u_j}\longrightarrow0,
\]
then $u_j\to u$ in capacity. Moreover, after passing to a subsequence, $P_\theta(\inf_{k\geq j}u_k)$ belongs to $\mathcal E_\varphi$ for every $j$.
\end{theorem}
\begin{proof}
Start with any subsequence and extract a further subsequence along which the integrals in the hypothesis are at most $2^{-j}$ and $u_j\to u$ almost everywhere for smooth volume. The envelope identity and Hartogs' lemma give a uniform bound $u_j,u\leq\varphi+C$. For $N\geq j$ let $p_{j,N}=P_\theta(u_j,\ldots,u_N)$. Partitioning the contact set as in the proof of the preceding lemma gives
\begin{equation}\label{eq:big-self-measure-roofs}
 \int_X|p_{j,N}-u|\,\mu_{p_{j,N}}\leq2^{1-j}.
\end{equation}
By \eqref{eq:big-self-measure-roofs}, \cref{prop:big-relative-first-moment}, applied with $w=u$ in its one-sided form, bounds $\sup_Xp_{j,N}$ from below, independently of $N$. Consequently the decreasing sequence has a limit
\[
 p_j=P_\theta\left(\inf_{k\geq j}u_k\right)
\]
that is not identically $-\infty$, and the same proposition puts $p_j$ in $\mathcal E_\varphi$. This proves the second assertion. Since $p_1\leq u_k$ for every $k$ and $\mathrm{e}^{u}\leq\mathrm{e}^{\sup_Xu}$, the preceding lemma, applied to the extracted sequence with $h=p_1$ and $\lambda=1$, shows that it converges to $u$ in capacity. Since the initial subsequence was arbitrary, the conclusion holds for the original sequence.
\end{proof}


\subsection{Higher moments in the minimal singularity type}
\label{sec:big-minimal-moments}

In this subsection the reference is the potential $V=V_\theta=P_\theta(0)$ with minimal singularities. It is an ordinary model, and we write $\mathcal E_V=\mathcal E_{\mathrm{env}}(X,\theta,V)$ and
\[
 \ell_V=\underline{\operatorname{vol}}_\theta(V)\geq\ell>0,
 \qquad M_V=\overline{\operatorname{vol}}_\theta(V)<\infty.
\]
The first inequality follows from volume monotonicity, since $\varphi\preceq V$, and the last from bounded mass. Thus $V$ satisfies the standing assumptions on $\varphi$, and the results above apply with $\varphi$ replaced by $V$. For a probability measure $\nu$, $A_\theta(\nu)$ denotes the supremum of $\int_X(-w)\,\mathrm{d}\nu$ over the $w\in\PSH(X,\theta)$ with $\sup_Xw=0$, as in \cref{subsec:psh-measure-smoothing}. We show that a first moment of $V-u$ against the measure $\mu_u$ of $u$ itself controls a higher moment of $V-u$ against every fixed $\nu$ with $A_\theta(\nu)<\infty$. An envelope with a concave obstacle converts the distribution of $V-u$ under $\mu_u$ into this stronger integrability estimate. We state the envelope lemma for the reference $\varphi$; here it is applied with $V$ in place of $\varphi$, and in \cref{sec:big-model-lp} with $\varphi$ itself.

\Needspace{20\baselineskip}
\begin{lemma}\label{lem:big-concave-envelope}\label{lem:big-model-concave-envelope}
Let $u\in\mathcal E_\varphi$ with $u\leq\varphi$, and let
$\chi\colon\mathbb R_{\leq0}\to\mathbb R_{\leq0}$ be increasing and concave, with $\chi(0)=0$. Suppose that $\chi$ is $C^{1,1}$ on compact intervals and that, for a finite constant $L\geq1$,
\[
 1\leq\chi'\leq L.
\]
Then the envelope
\[
 p=P_\theta(\varphi+\chi(u-\varphi))
\]
belongs to $\mathcal E_\varphi$, and
\[
 \mu_p\leq\mathbf1_{\{p=\varphi+\chi(u-\varphi)\}}
                 \chi'(u-\varphi)^n\mu_u.
\]
Here the obstacle and the contact set are defined off the pluripolar set where $u-\varphi$ is undefined, which all these measures ignore.
\end{lemma}
\begin{proof}
For $t\leq0$, the bounds on $\chi'$ give $Lt\leq\chi(t)\leq t$. Thus
\[
 Q_L(u,\varphi)\leq p\leq u\leq\varphi.
\]
The lower potential belongs to $\mathcal E_\varphi$ by \cref{lem:relative-stretch-closure}, so $p$ is not identically $-\infty$ and belongs to $\mathcal E_\varphi$ by upward stability. The obstacle is finite and quasi-continuous off the poles, so the contact theorem \cref{thm:big-general-envelope-contact} applies.

Let $\tau=\chi^{-1}$. It is convex, with $L^{-1}\leq\tau'\leq1$. We justify the transformation
\[
 w=\varphi+\tau(p-\varphi)
\]
without assuming that $\varphi$ is locally bounded on some open set. On a coordinate ball choose a smooth function $q$ with $\theta+\ddc q\geq0$. The functions
\[
 p_m=\max(p,q-m),\qquad
 \varphi_m=\max(\varphi,q-m)
\]
are locally bounded $\theta$-plurisubharmonic functions on the ball, and $p_m\leq\varphi_m$. \Cref{lem:bounded-convex-pair} applies to
$w_m=\varphi_m+\tau(p_m-\varphi_m)$ and gives
\[
 \theta_{w_m}\geq0,\qquad
 \mu_{w_m}\geq\tau'(p_m-\varphi_m)^n\mu_{p_m}.
\]

The function $(s,t)\mapsto t+\tau(s-t)$ on $\{s\leq t\}$ is nondecreasing in each variable. Consequently $w_m$ decreases as $m$ increases. Since $p\leq p_m\leq w_m\leq\varphi_m$, the limit is not identically $-\infty$ on the ball. Off the poles it equals $w$. The local limits therefore agree on overlaps after regularization and define a global $\theta$-plurisubharmonic function.

On $D_m=\{p>q-m\}$, neither $p$ nor $\varphi$ is truncated. Thus $w_m=w$ there, and plurifine locality transfers the product inequality to $D_m$. These sets exhaust the complement of the poles of $p$ in the ball. Since all products ignore the discarded pluripolar set,
\[
 \mu_w\geq\tau'(p-\varphi)^n\mu_p.
\]
The obstacle bound $p\leq\varphi+\chi(u-\varphi)$ gives $w\leq u$, with equality on the contact set of $p$. Ordered contact, the identity
$\tau'(p-\varphi)=1/\chi'(u-\varphi)$ on that set, and contact support of $\mu_p$ prove the assertion.
\end{proof}

\Needspace{11\baselineskip}
\begin{proposition}\label{prop:big-minimal-higher-moment}
Set $N=2n+1$ and $r=1+(4N)^{-1}$. For $u\in\PSH(X,\theta)$ with $u\sim V$ and $\sup_Xu=0$, put $B=\int_X(V-u)\,\mu_u$. There is a constant $C$, depending only on $n,\ell_V,M_V$, such that, for every non-pluripolar probability measure $\nu$ with $A=A_\theta(\nu)<\infty$,
\[
 \int_X(V-u)^r\,\mathrm{d}\nu
 \leq C(1+A)(1+B)^{1+1/N}
 .
\]
In particular, normalized potentials $u\sim V$ with uniformly bounded $B$ are uniformly integrable with respect to $\nu$.
\end{proposition}
\begin{proof}
From the distribution of $V-u$ under $\mu_u$ we build a concave weight adapted to $u$, apply the preceding lemma, and use $A$ to bound the resulting envelope against $\nu$. Throughout, $c$ and $C$ denote positive constants depending only on $n,\ell_V,M_V$.

The normalization gives $u\leq V$. Put $x=V-u\geq0$ and
\[
 T=\sup\{t\geq0:\mu_u(\{x>t\})>0\},
\]
with $T=0$ if this set is empty. Since $u\sim V$, $T$ is finite. Zero domination, applied to $u$ and $V-t$ for $t>T$, shows that $u\geq V-T$; thus $x\leq T$ off the pole set. If $T\leq2$ the claim is immediate. We may therefore suppose $T>2$; in particular $B>0$.

Fix a real number $s$ with $0<s<T-2$. For $t\geq0$, define
\[
 g_s(t)=
 \begin{cases}
  \bigl((1+t)^{3/2}\mu_u(\{x>t\})\bigr)^{-1},&0\leq t\leq T-s,\\
  (1+t)^{-3/2},&t>T-s,
 \end{cases}
 \quad F_s(t)=1+\int_0^t g_s(a)\,\mathrm{d}a,
\]
and set
\[
 W_s(t)=\int_0^t F_s(a)^{1/N}\,\mathrm{d}a,
 \qquad\chi_s(-t)=-W_s(t).
\]
The tail $\mu_u(\{x>t\})$ in the first denominator has a positive lower bound on $[0,T-s]$. Hence $F_s$ is bounded, nondecreasing and locally Lipschitz, with $F_s\geq1$. Since $u\leq V$ and, by \cref{prop:relative-elementary}, $u\in\mathcal E_V$, the function $\chi_s$ satisfies the hypotheses of \cref{lem:big-concave-envelope}, with $V$ in place of $\varphi$.

Tonelli's theorem gives
\begin{equation}\label{eq:big-weight-derivative-moment}
 \int_X F_s(x)\,\mu_u
 =\mu_u(X)+\int_0^Tg_s(t)\mu_u(\{x>t\})\,\mathrm{d}t
 \leq3\max(1,M_V),
\end{equation}
since $\int_0^\infty(1+t)^{-3/2}\,\mathrm{d}t=2$ and the tail is at most $M_V$ on the second interval.
Chebyshev's inequality gives $\mu_u(\{x>t\})\leq B/t$. Integrating first for $F_s$ and then for $W_s$, we obtain
\begin{equation}\label{eq:big-weight-growth}
 W_s(t)\geq c(1+B)^{-1/N}t^{\beta},
 \qquad1\leq t\leq T-s,
 \qquad\beta=1+\frac1{2N}.
\end{equation}
Indeed, $g_s(t)\geq cB^{-1}t^{-1/2}$ for $t\geq1$ in this interval, so $F_s(t)\geq c(1+B)^{-1}t^{1/2}$. For $t\geq2$ integrate this estimate over $[t/2,t]$; for $1\leq t\leq2$ use $W_s(t)\geq t$.

Let $p_s=P_\theta(V-W_s(x))\in\mathcal E_V$, the envelope of \cref{lem:big-concave-envelope} for $\chi=\chi_s$. The lemma, the bound $W_s(t)\leq tF_s(t)^{1/N}$ and the Cauchy--Schwarz inequality give, since $N=2n+1$,
\[
 \begin{split}
 \int_X(V-p_s)^{1/2}\,\mu_{p_s}
 &\leq\int_X x^{1/2}F_s(x)^{(n+1/2)/N}\,\mu_u\\
 &\leq B^{1/2}\left(\int_XF_s(x)\,\mu_u\right)^{1/2}
 \leq C B^{1/2}.
 \end{split}
\]
Write $a_s=\sup_Xp_s\leq0$. By the definition of $V$, $p_s-a_s\leq V$, so $V-p_s\geq-a_s$. Since $\mu_{p_s}(X)\geq\ell_V$, the last estimate yields $\ell_V\sqrt{-a_s}\leq C\sqrt B$, and hence $-a_s\leq CB$.

The definition of $A$ now gives
\[
 \int_XW_s(x)\,\mathrm{d}\nu
 \leq\int_X(V-p_s)\,\mathrm{d}\nu
 \leq A-a_s\leq A+CB.
\]
Combining this with \eqref{eq:big-weight-growth}, and then letting $s\downarrow0$ for each fixed $1\leq t<T$, we obtain
\[
 \nu(\{x>t\})\leq C(1+B)^{1/N}(A+B)t^{-\beta}.
\]
There is no tail for $t\geq T$. The layer-cake formula therefore gives
\[
 \int_Xx^r\,\mathrm{d}\nu
 =r\int_0^\infty t^{r-1}\nu(\{x>t\})\,\mathrm{d}t
 \leq1+C(1+B)^{1/N}(A+B),
\]
because $1<r<\beta$. This implies the stated bound.

Finally $-u=-V+(V-u)$, where $-V\in L^1(\nu)$ because $A<\infty$. The first summand is fixed and integrable, and the second has a uniform $L^r(\nu)$ bound. This proves uniform integrability; no $L^r$ bound for $V$ itself is claimed.
\end{proof}


\section{General measures in the minimal class}
\label{sec:big-measure-smoothing}

We now solve the equations in $\mathcal E_V$. For a measure $\nu$ with $A_\theta(\nu)<\infty$, the coordinate convolution of \cref{subsec:psh-measure-smoothing} gives approximating equations with smooth densities. The integral estimates for this smoothing, combined with the higher-moment bound, give convergence of the solutions against the original measure. The resulting self-measure error yields convergence in capacity, and a first-moment tail estimate identifies the limit of the Monge--Amp\`ere measures.

\Needspace{18\baselineskip}
\begin{proposition}\label{prop:big-smoothed-measure-continuity}
Let $\nu$ be a non-pluripolar probability measure with $A_\theta(\nu)<\infty$, and let $\nu_j$ be its smooth approximations defined by \eqref{eq:measure-smoothing-duality}. Fix a constant $C>0$ and let $v_j,v\in\PSH(X,\theta)$ satisfy
\[
 v_j\sim V,\qquad \sup_Xv_j=0,\qquad
 \mu_{v_j}\leq C\nu_j,
 \qquad v_j\longrightarrow v\text{ in }L^1(X).
\]
Then $v\in\mathcal E_V$, the convergence also holds in $L^1(\nu)$ and in capacity, and
\[
 \mu_{v_j}\rightharpoonup\mu_v.
\]
\end{proposition}
\begin{proof}
The limit $v$ is not identically $-\infty$ and satisfies $\sup_Xv=0$, and $v_j,v\leq V\leq0$. By \cref{lem:measure-smoothing-integrals}, $\sup_jA_\theta(\nu_j)<\infty$, which gives the uniform estimate
\begin{equation}\label{eq:smoothed-first-moment}
 \int_X(-v_j)\,\mu_{v_j}
 \leq C\int_X(-v_j)\,\mathrm{d}\nu_j\leq C'.
\end{equation}
In particular the relative first moments $\int_X(V-v_j)\mu_{v_j}$ are bounded, and \cref{prop:big-relative-first-moment}, with fixed reference $w=V$, puts $v$ in $\mathcal E_V$.

By \cref{prop:big-minimal-higher-moment}, applied to the fixed measure $\nu$, the $v_j$ are uniformly integrable with respect to $\nu$. For each fixed $T$, the bounded functions $\max(v_j,-T)$ have a common lower curvature bound in a Hermitian background $\Omega\geq\theta$ and converge in $L^1$ for smooth volume to $\max(v,-T)$; hence \cref{lem:fixed-npp-integrals} gives their convergence in $L^1(\nu)$. These truncations need not be $\theta$-plurisubharmonic; they are used only in the larger Hermitian background. Uniform integrability, followed by $T\to\infty$, gives $v_j\to v$ in $L^1(\nu)$, and the limit is integrable because $A_\theta(\nu)<\infty$.

The smoothing lemma, in the form \eqref{eq:convolution-moving-integrals}, now yields
\[
 \int_X|v_j-v|\,\mu_{v_j}
 \leq C\int_X|v_j-v|\,\mathrm{d}\nu_j\longrightarrow0.
\]
Therefore the self-measure convergence theorem \cref{thm:big-relative-self-measure-convergence} gives convergence in capacity. Finally, \eqref{eq:smoothed-first-moment} gives the sublevel estimate
\[
 \sup_j\mu_{v_j}(\{v_j\leq-T\})\leq C'/T,
\]
and the masses of the $\mu_{v_j}$ are bounded by $M_V$. The tail criterion \cref{thm:general-background-damping} therefore gives the asserted weak convergence.
\end{proof}

\Needspace{36\baselineskip}
\subsection{Exponential and normalized equations}
\label{sec:big-minimal-general-measures}

The smoothing argument solves both equations for measures with finite $A_\theta$. A bounded carrier representation reduces an arbitrary finite non-pluripolar measure to increasing bounded densities, whose solutions have measures dominated by one fixed finite measure; \cref{thm:big-relative-fixed-continuity} then passes to the limit.

\begin{theorem}\label{thm:big-minimal-general-measures}
For every finite non-pluripolar measure $\mu$ with $\mu(X)>0$, the following assertions hold.
\begin{enumerate}
\item For every $\lambda>0$ there is a unique $u\in\mathcal E_V$ such that
\[
 \mathrm{MA}_\theta(u)=\mathrm{e}^{\lambda u}\mu.
\]
\item There are $v\in\mathcal E_V$ with $\sup_Xv=0$ and a unique scalar $c>0$ such that
\[
 \mathrm{MA}_\theta(v)=c\mu.
\]
\end{enumerate}
\end{theorem}
\begin{proof}
First suppose that $\mu$ is a probability measure with $A_\theta(\mu)<\infty$, and let $\mu_j$ be its smooth probability approximations. The theorems for the minimal ordinary type, \cref{thm:big-minimal-exponential,thm:big-minimal-normalized}, apply to their smooth, nonzero densities.

For the normalized equation, they give $v_j\sim V$ with $\sup_Xv_j=0$ and $\mu_{v_j}=c_j\mu_j$. Since $\mu_j$ is a probability measure, the mass bounds $\ell_V\leq\mu_{v_j}(X)\leq M_V$ give $\ell_V\leq c_j\leq M_V$. After extraction, $c_j\to c>0$ and $v_j\to v$ in $L^1(X)$. The preceding proposition applies, so $v\in\mathcal E_V$ and $\mu_{v_j}\rightharpoonup\mu_v$. The right-hand sides converge weakly to $c\mu$, which proves the equation.

For the positive-exponent equation, take $u_j\sim V$ satisfying $\mu_{u_j}=\mathrm{e}^{\lambda u_j}\mu_j$, and write $a_j=\sup_Xu_j$ and $v_j=u_j-a_j$. The mass bounds and Jensen's inequality give
\[
 \ell_V\leq \mathrm{e}^{\lambda a_j},\qquad
 M_V\geq \mathrm{e}^{\lambda a_j}\int_X\mathrm{e}^{\lambda v_j}\,\mathrm{d}\mu_j
       \geq \mathrm{e}^{\lambda a_j-\lambda A_\theta(\mu_j)},
\]
and $\sup_jA_\theta(\mu_j)<\infty$ by \cref{lem:measure-smoothing-integrals}. Thus $a_j$ is bounded above and below. After extraction, $a_j\to a$, $v_j\to v$ in $L^1(X)$, and $\mu_{v_j}\leq C\mu_j$. The preceding proposition gives $v\in\mathcal E_V$, convergence in $L^1(\mu)$ and weak convergence of the Monge--Amp\`ere measures. Set $u=v+a$. By the smoothing lemma, the right-hand sides converge weakly to $\mathrm{e}^{\lambda u}\mu$. This proves the equation, and $u\in\mathcal E_V$.

The same conclusions hold for every nonzero finite measure whose normalization to a probability measure has finite $A_\theta$, since scaling the measure only changes the positive scalar, or translates the solution of the exponential equation by a constant.

Now let $\mu$ be an arbitrary finite non-pluripolar measure of positive mass. The bounded-carrier decomposition \cref{lem:bounded-measure-carrier}, applied to $\omega$, writes
\[
 \mu=f\nu,\qquad \nu=\mathrm{MA}_\omega(\psi),\qquad
 \psi\in\PSH(X,\omega)\cap L^\infty(X),\qquad
 0\leq f\in L^1(\nu).
\]
For $\mu_k=\min(f,k)\nu$, we have $\mu_k(X)>0$. By \cref{lem:bounded-carrier-integrability}, the normalization of $\mu_k$ has finite $A_\theta$ for each fixed $k$.

For the normalized equation, the preceding case gives $v_k\in\mathcal E_V$ with $\sup_Xv_k=0$ and $\mu_{v_k}=c_k\mu_k$, where
\[
 \frac{\ell_V}{\mu(X)}\leq c_k\leq\frac{M_V}{\mu_1(X)}.
\]
After extraction, $c_k\to c>0$ and $v_k$ converges in $L^1$ to some $v\not\equiv-\infty$. All the measures $\mu_{v_k}$ are bounded by the fixed finite non-pluripolar measure $M_V\mu/\mu_1(X)$. The fixed-dominator theorem \cref{thm:big-relative-fixed-continuity} gives $v\in\mathcal E_V$ and weak convergence to $\mu_v$. Since $c_k\mu_k\to c\mu$ in total variation, $\mu_v=c\mu$.

For the positive-exponent equation, let $u_k\in\mathcal E_V$ solve $\mu_{u_k}=\mathrm{e}^{\lambda u_k}\mu_k$. The argument that follows the construction of the approximating solutions in the positive-exponent part of the proof of \cref{thm:general-measure-minimal} applies, with \cref{cor:big-relative-orders,thm:big-relative-fixed-continuity} and $\ell_V$ in place of \cref{cor:relative-extended-exponential-order,thm:relative-fixed-measure-continuity} and $v_-(\theta)$. It shows that the $u_k$ decrease to some $u\in\mathcal E_V$ with $\mu_u=\mathrm{e}^{\lambda u}\mu$.

Finally, \cref{cor:big-relative-orders} with reference $V$, applied in both directions, gives uniqueness of the exponential solution and of the positive scalar.
\end{proof}

\Needspace{23\baselineskip}
Positive lower mass of a Hermitian form gives positive lower volume to the minimal reference of every smooth big background.

\begin{corollary}\label{cor:metric-pvp-minimal-volume}
Suppose that $v_-(\omega)>0$. Then $\underline{\operatorname{vol}}_\theta(V_\theta)>0$ for every smooth big real $(1,1)$-form $\theta$ on $X$; hence the preceding theorem applies to every such $\theta$.
\end{corollary}
\begin{proof}
Choose a strict barrier $\rho\in\PSH(X,\theta)$, so that $\theta_\rho\geq a\omega$ for some $a>0$. Then $\rho\in\PSH(X,\theta-a\omega)$. Volume addition (\cref{prop:volume-addition}), followed by monotonicity in the singularities (\cref{thm:volume-monotonicity}), gives
\[
 0<a^n v_-(\omega)
 =\underline{\operatorname{vol}}_{a\omega}(0)
 \leq\underline{\operatorname{vol}}_\theta(\rho)
 \leq\underline{\operatorname{vol}}_\theta(V_\theta).
\]
In the addition step the big background is $a\omega$, and the added positive current is $(\theta-a\omega)_\rho$. Bounded mass gives finiteness of the upper volume.
\end{proof}


\section{Weighted measure addition}
\label{sec:big-weighted-limits}

Given $u,v\in\mathcal E_\varphi$ and $\lambda>0$, we look for $a\in\mathcal E_\varphi$ with $\mu_a=\mathrm{e}^{\lambda(a-u)}\mu_u+\mathrm{e}^{\lambda(a-v)}\mu_v$. This weighted addition is used in \cref{sec:big-model-general-measures} to enlarge the weighted measure of a supersolution. The measure $\mathrm{e}^{-\lambda u}\mu_u+\mathrm{e}^{-\lambda v}\mu_v$ can have infinite mass, so we approximate the inputs while keeping a common lower bound in $\mathcal E_\varphi$ for the approximate solutions. This bound gives convergence in capacity, even when the approximate solutions have different ordinary types, and localizes the limiting equation away from deep sublevel sets. Exponentially weighted measures are defined off the pole sets and assigned zero mass there.

\Needspace{13\baselineskip}
\begin{lemma}\label{lem:big-weighted-common-bound}
Let $a_j\in\PSH(X,\theta)$ have a common upper bound and converge in $L^1(X)$ to $a\in\PSH(X,\theta)$ with $a\preceq\varphi$. Suppose there is $h\in\mathcal E_\varphi$ such that $h\leq a_j$ for every $j$. If, for some $\lambda>0$,
\[
 \int_X|a_j-a|\mathrm{e}^{\lambda a}\,\mu_{a_j}\longrightarrow0,
\]
then $a\in\mathcal E_\varphi$ and $a_j\to a$ in capacity.
\end{lemma}
\begin{proof}
By \cref{lem:qpsh-almost-everywhere-order}, the lower bound passes to the limit, so $h\leq a\preceq\varphi$, and upward stability gives $a\in\mathcal E_\varphi$. For the convergence in capacity we follow the proof of \cref{thm:big-relative-self-measure-convergence}, with $a_j,a$ in place of $u_j,u$, the weights $\min(1,|p_{j,N}-a|)\mathrm{e}^{\lambda a}$ in place of $\min(1,|p_{j,N}-u|)$, and the fixed lower bound $h$ in place of $p_1$. Partitioning contact sets gives, along a subsequence with integrals at most $2^{-j}$ and almost-everywhere convergence,
\[
 \int_X\min(1,|p_{j,N}-a|)\mathrm{e}^{\lambda a}\,\mu_{p_{j,N}}
 \leq\sum_{k=j}^N\int_X|a_k-a|\mathrm{e}^{\lambda a}\,\mu_{a_k}
 \leq2^{1-j}
\]
for the finite rooftops $p_{j,N}=P_\theta(a_j,\ldots,a_N)$, which are not identically $-\infty$ because $h$ is a competitor. The first-moment criterion is not needed: the decreasing limits $p_j$ satisfy $h\leq p_j\leq a$, so the $p_j$ and their increasing regularized limit $p$ lie in $\mathcal E_\varphi$, although the finite rooftops need not. The weights converge in capacity: discard deep sublevel sets of $h$ and $a$, and use continuity on the remaining finite ranges. Since the exponential is positive off a pluripolar set, which $\mu_p$ does not charge, the limiting identity $\int_X\min(1,|p-a|)\mathrm{e}^{\lambda a}\,\mu_p=0$ again gives $\mu_p(\{p<a\})=0$, and the rest of that proof applies.
\end{proof}

The next lemma treats exponent one; the general exponent in \cref{thm:big-weighted-addition} follows by rescaling the background and the potentials.

\Needspace{20\baselineskip}
\begin{lemma}\label{lem:big-truncated-addition-limit}
Let $u,v,\sigma\in\PSH(X,\theta)$ be nonpositive, with $u,v\preceq\varphi$, and let $t_j\uparrow\infty$. Each of the sequences $u_j$, $v_j$ is either constant, equal to $u$, respectively $v$, for every $j$, or given by the truncations
\[
 u_j=\max(u,\sigma-t_j),\qquad
 v_j=\max(v,\sigma-t_j).
\]
Fix $h\in\mathcal E_\varphi$ and let $a_j\in\PSH(X,\theta)$ satisfy
\[
 \mu_{a_j}=\mathrm{e}^{a_j-u_j}\mu_{u_j}+\mathrm{e}^{a_j-v_j}\mu_{v_j},
 \qquad h\leq a_j\leq\min(u_j,v_j).
\]
If $a_j\to a$ in $L^1(X)$, then $a\in\mathcal E_\varphi$, the convergence holds in capacity, and
\[
 a\leq\min(u,v),\qquad
 \mu_a=\mathrm{e}^{a-u}\mu_u+\mathrm{e}^{a-v}\mu_v.
\]
\end{lemma}
\begin{proof}
Since $u_j$ and $v_j$ decrease to $u$ and $v$, the bounds give $h\leq a\leq\min(u,v)\preceq\varphi$, so $a\in\mathcal E_\varphi$. After discarding finitely many indices we may assume $t_j\geq0$, so all potentials $a_j,a$ are nonpositive. Let $E_j$ be the intersection of the sets $\{u>\sigma-t_j\}$ and $\{v>\sigma-t_j\}$, taken only for those inputs that are truncated; a constant sequence imposes no condition. On $E_j$, locality gives $\mu_{u_j}=\mu_u$ and $\mu_{v_j}=\mu_v$. Since $a,a_j\leq u_j,v_j$, we have
\[
 \mathrm{e}^a\mu_{a_j}
 \leq \mathrm{e}^{(a+a_j)/2}(\mu_{u_j}+\mu_{v_j}).
\]
If $x\notin E_j$, then one truncated input and its truncation are at most $-t_j$ at $x$, so $a_j(x),a(x)\leq-t_j$. Moreover $|a_j-a|\mathrm{e}^{(a+a_j)/4}$ is bounded by an absolute constant, since $a_j,a\leq0$. Thus
\begin{equation}\label{eq:big-addition-weighted-error}
 \int_X|a_j-a|\mathrm{e}^a\,\mu_{a_j}
 \leq\int_X|a_j-a|\mathrm{e}^{(a+a_j)/2}(\mu_u+\mu_v)
       +C \mathrm{e}^{-t_j/2}.
\end{equation}
The last term uses the uniform bound for the masses of the $u_j$ and $v_j$ given by \cref{prop:bmp-finite}. The fixed measure $\mu_u+\mu_v$ is finite by bounded mass. \Cref{lem:fixed-npp-integrals}, applied to $\mathrm{e}^{a_j}$ and $\mathrm{e}^a$, shows that $a_j\to a$ in measure for $\mu_u+\mu_v$, since $a$ is finite almost everywhere for this non-pluripolar measure. The integrand in the first term of \eqref{eq:big-addition-weighted-error} is uniformly bounded and tends to zero in measure, so its integral tends to zero. \Cref{lem:big-weighted-common-bound} now gives convergence in capacity. When a subsequence is taken, the corresponding truncation levels $t_j$ are kept.

To identify the equation, fix $T>0$ and choose a continuous nondecreasing function $g_T\colon[-\infty,0]\to[0,1]$ that vanishes on $[-\infty,-2T]$ and equals one on $[-T,0]$, such that $g_T\uparrow1$ at finite points as $T\to\infty$. Put $\chi_T=g_T(h)$. Since $h\leq a\leq u,v$, both $u$ and $v$ exceed $-2T$ on $\{\chi_T>0\}$. Therefore, once $t_j>2T$, every truncated input agrees there with its limit, and
\[
 \chi_T\mu_{a_j}
 =\chi_T\mathrm{e}^{a_j}(\mathrm{e}^{-u}\mu_u+\mathrm{e}^{-v}\mu_v).
\]
\Needspace{7\baselineskip}
Choose a Hermitian form $\Omega\geq\theta$. On the left-hand side, replace $a_j$ by $\max(a_j,-2T)$, which does not change $\chi_T\mu_{a_j}$ because $\chi_T$ vanishes where $a_j\leq-2T$, and expand the bounded product as
\[
 (\theta+\ddc\max(a_j,-2T))^n
 =\sum_{k=0}^n\binom nk
   (\Omega+\ddc\max(a_j,-2T))^k\wedge(\theta-\Omega)^{n-k}.
\]
Bounded capacity continuity with the fixed quasi-continuous weight $\chi_T$ therefore gives convergence to $\chi_T\mu_a$, after testing against a continuous function. The uncut product may be signed, but its terms have uniformly capacity-controlled total variations.

On the right-hand side, $\chi_T \mathrm{e}^{-u}$ and $\chi_T \mathrm{e}^{-v}$ are bounded by $\mathrm{e}^{2T}$. Fixed-measure exponential convergence therefore gives convergence in total variation to
$\chi_T\mathrm{e}^a(\mathrm{e}^{-u}\mu_u+\mathrm{e}^{-v}\mu_v)$. Let $T\to\infty$. The limiting right-hand side is at most $\mu_u+\mu_v$, because $a\leq u,v$. Both sides ignore the pole set of $h$, so the cutoff identities prove the global equation.
\end{proof}


\subsection{Ordinary types and bounded carriers}
\label{sec:big-weighted-addition}

We first solve the weighted addition equation when one input has the ordinary type of $\varphi$, or a measure dominated by the Monge--Amp\`ere measure of a bounded potential. In both cases first moments control the rooftops of the approximating solutions, and these rooftops give the common lower bound required by \cref{lem:big-truncated-addition-limit}. The solution lies below both inputs, so its measure is finite.

\Needspace{14\baselineskip}
\begin{theorem}\label{thm:big-weighted-addition}
Let $u,v\in\mathcal E_\varphi$ and $\lambda>0$. Suppose either $v\sim\varphi$, or there are a Hermitian form $\omega$, a bounded $\psi\in\PSH(X,\omega)$ and a finite constant $A>0$ such that
\[
 \mu_v\leq A\,\mathrm{MA}_\omega(\psi).
\]
Then there is a unique $a\in\mathcal E_\varphi$ satisfying
\[
 \mu_a=\mathrm{e}^{\lambda(a-u)}\mu_u+\mathrm{e}^{\lambda(a-v)}\mu_v.
\]
It satisfies $a\leq\min(u,v)$, and hence $\mu_a\leq\mu_u+\mu_v$. If both $u,v\sim\varphi$, then also $a\sim\varphi$.
\end{theorem}
\begin{proof}
The proof has four steps: two inputs of minimal singularity type, two inputs of the ordinary type of $\varphi$, one input of that type, and finally an input with a bounded carrier measure.

It suffices to prove the result for $\lambda=1$. Indeed, replace $(\theta,\varphi,u,v)$ by $(\lambda\theta,\lambda\varphi,\lambda u,\lambda v)$: envelopes and relative classes scale with the potentials, volumes and Monge--Amp\`ere measures scale by $\lambda^n$, and a bounded-carrier domination remains one after changing its constant. Dividing the resulting solution by $\lambda$ gives the desired equation. Similarly, a common downward translation of $u$, $v$ and the solution allows us to assume $u,v\leq\varphi\leq0$; by the envelope identity, these inequalities hold after a sufficiently large common translation. We restore the translation at the end.

\emph{Step 1: inputs of minimal singularity type.}
First take $\varphi=V$ and $u,v\sim V$. Truncating only the weights, define the finite nonzero measures
\[
 \tau_j=\mathrm{e}^{-\max(u,-j)}\mu_u+\mathrm{e}^{-\max(v,-j)}\mu_v.
\]
The functions $\max(u,-j)$ need not be $\theta$-plurisubharmonic; they occur only as weights. \Cref{thm:big-minimal-general-measures} gives $a_j\in\mathcal E_V$ with $\mu_{a_j}=\mathrm{e}^{a_j}\tau_j$.

Choose $D\geq\tfrac12\|u-v\|_\infty+n\log2$, where the difference is taken off the common poles, and put $w=(u+v)/2-D$. Positivity of mixed products gives
\[
 \mu_w\geq2^{-n}(\mu_u+\mu_v)\geq \mathrm{e}^w\tau_j.
\]
Exponential order in $\mathcal E_V$ (\cref{cor:big-relative-orders}) gives $a_j\geq w$. Since $\tau_j$ increases with $j$, the same order shows that $a_j$ decreases. The envelope identity gives $a_1\leq V+\sup_Xa_1$, so $a_1\sim V$ by the lower bound $w$. Consequently
\[
 \mu_{a_j}
 \leq \mathrm{e}^{a_1-u}\mu_u+\mathrm{e}^{a_1-v}\mu_v
 \leq C(\mu_u+\mu_v).
\]
\Needspace{8\baselineskip}
The decreasing limit $a\geq w$ is not identically $-\infty$ and has the ordinary type of $V$. Fixed-dominator continuity (\cref{thm:big-relative-fixed-continuity}) gives weak convergence of the left-hand sides, and the right-hand sides converge in total variation by dominated convergence with respect to $\mu_u+\mu_v$. Therefore
\[
 \mu_a=\mathrm{e}^{a-u}\mu_u+\mathrm{e}^{a-v}\mu_v.
\]
Comparing this equation with each input measure, exponential order now gives $a\leq u,v$. This comparison uses the limiting equation; the truncated weights alone would not give it.

\emph{Step 2: two inputs of the ordinary type of $\varphi$.}
Suppose $u,v\sim\varphi$, and put $u_j=\max(u,V-j)$ and $v_j=\max(v,V-j)$. Step 1 gives solutions $a_j\sim V$ with
\[
 \mu_{a_j}=\mathrm{e}^{a_j-u_j}\mu_{u_j}+\mathrm{e}^{a_j-v_j}\mu_{v_j},
 \qquad a_j\leq\min(u_j,v_j).
\]
The differences $u_j-v_j$ satisfy the same uniform bound as $u-v$. With $D$ as above, $w_j=(u_j+v_j)/2-D$ satisfies the reverse inequality for this equation, and exponential order gives
\[
 (u+v)/2-D\leq w_j\leq a_j\leq\min(u_j,v_j).
\]
The fixed lower bound $w=(u+v)/2-D$ has the ordinary type of $\varphi$, and hence lies in $\mathcal E_\varphi$. Extract an $L^1$ limit $a\not\equiv-\infty$; it satisfies $w\leq a\leq\min(u,v)$, so $a\sim\varphi$. \Cref{lem:big-truncated-addition-limit}, with $\sigma=V$ and lower bound $h=w$, proves convergence in capacity and the equation. The lemma applies although the approximating solutions have the type of $V$ and the limit has the type of $\varphi$.

\emph{Step 3: one input of the ordinary type of $\varphi$.}
Now let $u\in\mathcal E_\varphi$ and $v\sim\varphi$, and set $u_j=\max(u,\varphi-j)$. Step 2 gives $a_j\sim\varphi$ with
\[
 \mu_{a_j}=\mathrm{e}^{a_j-u_j}\mu_{u_j}+\mathrm{e}^{a_j-v}\mu_v,
 \qquad a_j\leq\min(u_j,v).
\]
We construct a common lower bound in $\mathcal E_\varphi$. Put $\delta=n\log2>0$ and
\[
 s_j=Q_2(a_j,u_j)\in\mathcal E_\varphi,
 \qquad C_j=\{s_j+u_j=2a_j\}.
\]
\Needspace{6\baselineskip}
By stretching and contact, $s_j\leq a_j$ and
\begin{equation}\label{eq:big-addition-stretch-contact}
 2^{-n}\mu_{s_j}+2^{-n}\mathbf1_{C_j}\mu_{u_j}
 \leq\mathbf1_{C_j}\mu_{a_j}.
\end{equation}
Put $D_j=\{s_j\leq u-2\delta\}$. On $D_j\cap C_j$ we have $a_j=(s_j+u_j)/2\leq u_j-\delta$, because $u\leq u_j$. Since $\mathrm{e}^{-\delta}=2^{-n}$, inserting the equation for $\mu_{a_j}$ into \eqref{eq:big-addition-stretch-contact} shows that the $\mu_{u_j}$ terms cancel there. Also $v\geq\varphi-C\geq u_j-C$ for a fixed constant $C$. Hence
\begin{equation}\label{eq:big-addition-ordinary-tail}
 \mathbf1_{D_j}\mu_{s_j}
 \leq2^n \mathrm{e}^C\mathbf1_{D_j\cap C_j}
                \mathrm{e}^{(s_j-u_j)/2}\mu_v.
\end{equation}
All measures being subtracted here are finite.

\Needspace{6\baselineskip}
For $N\geq1$ form the finite rooftops
\[
 R_N=P_\theta(s_1,\ldots,s_N),\qquad q_N=P_\theta(u,R_N).
\]
They belong to $\mathcal E_\varphi$ by the finite-rooftop proposition. On the set $\{R_N\leq u-2\delta\}$, partition the contact set according to the first index $i$ with $R_N=s_i$. Since $u_i\geq u_N$, \eqref{eq:big-addition-ordinary-tail} gives
\[
 \mathbf1_{\{R_N\leq u-2\delta\}}\mu_{R_N}
 \leq2^n \mathrm{e}^C \mathrm{e}^{(R_N-u_N)/2}\mu_v.
\]
\Needspace{7\baselineskip}
The contact pieces are disjoint, so there is no factor $N$. On $\{q_N\leq u-2\delta\}$ only the contact part $q_N=R_N$ contributes. Since $u\leq u_N$ and $t \mathrm{e}^{-t/2}\leq2/\mathrm{e}$ for $t\geq0$, we obtain
\[
 \int_X(u-q_N)\,\mu_{q_N}
 \leq 2\delta M+2^{n+1}\mathrm{e}^{C-1}\mu_v(X).
\]
The first term bounds the contribution of the complementary set, including any contact with $u$. The first-moment criterion \cref{prop:big-relative-first-moment}, with fixed reference $u$, implies that the decreasing sequence $q_N$ has a limit $h\in\mathcal E_\varphi$, not identically $-\infty$. For each fixed $j$ and every $N\geq j$, we have $q_N\leq s_j\leq a_j$, so $h\leq a_j$ for every $j$.

The $a_j$ now have a subsequence converging in $L^1$ to a limit not identically $-\infty$, bounded above by $u$ and $v$. \Cref{lem:big-truncated-addition-limit}, with reference $\varphi$, fixed second input $v$ and common lower bound $h$, proves the required equation in $\mathcal E_\varphi$.

\emph{Step 4: a bounded carrier measure for the second input.}
Finally let $u,v\in\mathcal E_\varphi$ and $\mu_v\leq A\nu$, with $\nu=\mathrm{MA}_\omega(\psi)$ and $\psi$ bounded. Again put $u_j=\max(u,\varphi-j)$. By Step 3, applied with the inputs interchanged, there are $a_j\in\mathcal E_\varphi$ satisfying the same equations and bounds as in Step 3.

Let $r=P_\theta(u,v)\in\mathcal E_\varphi$. On $\{a_j<r\}$ one has $r\leq u_j,v$, so the inequality $t \mathrm{e}^{-t}\leq1/\mathrm{e}$ gives
\[
 \int_X(r-a_j)_+\,\mu_{a_j}
 \leq \mathrm{e}^{-1}\bigl(\mu_{u_j}(X)+\mu_v(X)\bigr)\leq2M/\mathrm{e}.
\]
\Needspace{7\baselineskip}
Since $a_j\leq\varphi$, the one-sided first-moment criterion gives a uniform lower bound for $\sup_Xa_j$. By \cref{lem:bounded-carrier-integrability},
\[
 \sup_{\substack{z\in\PSH(X,\theta)\\\sup_Xz=0}}
       \int_X(-z)\,\mathrm{d}\nu<\infty.
\]
\Needspace{6\baselineskip}
Applying this to $a_j-\sup_Xa_j$ and using the lower bound for the suprema, we obtain
\begin{equation}\label{eq:big-addition-carrier-moment}
 \sup_j\int_X(u_j-a_j)\,\mu_v
 \leq\sup_j\int_X(-a_j)\,\mu_v<\infty.
\end{equation}

Keep $s_j=Q_2(a_j,u_j)$, $C_j$ and $D_j$ from Step 3. The same cancellation in \eqref{eq:big-addition-stretch-contact}, now using only $a_j\leq v$, gives
\[
 \mathbf1_{D_j}\mu_{s_j}
 \leq2^n\mathbf1_{D_j\cap C_j}\mathrm{e}^{a_j-v}\mu_v
 \leq2^n\mu_v.
\]
\Needspace{6\baselineskip}
Take one additional rooftop,
\[
 q_j=P_\theta(s_j,u-2\delta)\in\mathcal E_\varphi.
\]
On its contact set with $u-2\delta$ its measure is at most $\mu_u$; on its other contact part, where $q_j=s_j\leq u-2\delta$, the preceding inequality applies. Thus
\begin{equation}\label{eq:big-addition-fixed-roof-measure}
 \mu_{q_j}\leq\mu_u+2^n\mu_v.
\end{equation}
Moreover the integrand $u-2\delta-q_j$ vanishes on the first contact part. On the other part, the measure is carried by $C_j$, where $u_j-s_j=2(u_j-a_j)$. Therefore
\[
 \int_X(u-2\delta-q_j)\,\mu_{q_j}
 \leq2^{n+1}\int_X(u_j-a_j)\,\mu_v,
\]
which is uniformly bounded by \eqref{eq:big-addition-carrier-moment}. The first-moment criterion with reference $u-2\delta$ prevents collapse of the $q_j$. Extract $L^1$ limits of the $a_j$ and of the $q_j$, not identically $-\infty$. The fixed-dominator theorem applies to \eqref{eq:big-addition-fixed-roof-measure} and, after a further extraction, gives
\[
 h=P_\theta\left(\inf_jq_j\right)\in\mathcal E_\varphi.
\]
Since $q_j\leq s_j\leq a_j$, this is a common lower bound for the extracted solutions, and \cref{lem:big-truncated-addition-limit} gives the equation. Its truncation levels are the original indices of the extracted subsequence.

In every step the solution lies below $u$ and $v$, so its measure is bounded by $\mu_u+\mu_v$. Two solutions have the same exponentially weighted measure, so \cref{cor:big-relative-orders}, applied in both directions, gives uniqueness, even when that weighted measure has infinite mass. Restoring the common translation and the exponent finishes the proof.
\end{proof}


\Needspace{10\baselineskip}
\section{Relative projection and supersolutions}
\label{sec:big-relative-projection}

To pass from the minimal class to the reference $\varphi$, we project a solution $z\in\mathcal E_V$ onto $\mathcal E_\varphi$. The projection is the increasing limit $P_\theta[\varphi](z)$ of the rooftops $P_\theta(\varphi+C,z)$ as $C\to\infty$, in the notation of \cref{lem:relative-obstacle-contact}. The small stretches of \cref{lem:positive-volume-stretch} first show that the rooftop $P_\theta(\varphi,z)$ is not identically $-\infty$ and has the full upper volume of $\varphi$. Contact support then bounds the measure of the projection by the measure of $z$ on the set where the two coincide. The argument uses only that the reference of $z$ lies above $\varphi$, so we formulate it for an arbitrary ordinary model above $\varphi$; this will also allow the reference to vary.

Let $\eta\in\PSH(X,\theta)$ satisfy $P_\theta[\eta]=\eta\geq\varphi$, and write $\mathcal E_\eta=\mathcal E_{\mathrm{env}}(X,\theta,\eta)$.

\begin{proposition}\label{prop:big-cross-class-roof}
For every $z\in\mathcal E_\eta$, the rooftop $P_\theta(\varphi,z)$ is not identically $-\infty$ and belongs to $\mathcal E_\varphi$.
\end{proposition}
\Needspace{7\baselineskip}
\begin{proof}
First suppose $z\leq\eta$; recall also that $\varphi\leq\eta\leq0$. For $b>1$ sufficiently close to one, \cref{lem:positive-volume-stretch} gives
\[
 h_b=Q_b(\varphi,\eta)\leq\varphi,
\]
which is not identically $-\infty$. Set $a=1-b^{-1}$. The stretching lemma \cref{lem:relative-stretch-closure}, now applied in $\mathcal E_\eta$, gives $z_{1/a}=Q_{1/a}(z,\eta)\in\mathcal E_\eta$. The potential
\[
 k_a=a z_{1/a}+(1-a)h_b
\]
lies below both $\varphi$ and $z$. Indeed, $z_{1/a}\leq\eta$ and $(1-a)h_b+a\eta\leq\varphi$ give the first bound, while $h_b\leq\eta$ and $a z_{1/a}+(1-a)\eta\leq z$ give the second. Thus $p=P_\theta(\varphi,z)$ is not identically $-\infty$. Volume addition and monotonicity give
\[
 (1-a)^n\overline{\operatorname{vol}}_\theta(h_b)
 \leq\overline{\operatorname{vol}}_\theta(k_a)
 \leq\overline{\operatorname{vol}}_\theta(p)
 \leq\overline{\operatorname{vol}}_\theta(\varphi).
\]
As $b\downarrow1$, we have $a\downarrow0$, and the first term tends to the last by \cref{lem:positive-volume-stretch}. The upper volumes of $p$ and $\varphi$ are therefore equal. Since $p\leq\varphi$, \cref{thm:relative-equivalence} gives $p\in\mathcal E_\varphi$.

\Needspace{6\baselineskip}
For general $z$, the envelope identity in $\mathcal E_\eta$ gives $z-\sup_Xz\leq\eta$. Set $S=\max(\sup_Xz,0)$. The preceding case gives $P_\theta(\varphi,z-S)\in\mathcal E_\varphi$, and
\[
 P_\theta(\varphi,z-S)\leq P_\theta(\varphi,z)\leq\varphi.
\]
Upward stability finishes the proof.
\end{proof}

\Needspace{12\baselineskip}
\begin{proposition}\label{prop:big-relative-projection}\index{Envelope!relative projection}
Let $z\in\mathcal E_\eta$ and let $Q=P_\theta[\varphi](z)=\bigl(\lim_{C\to\infty}P_\theta(\varphi+C,z)\bigr)^*$ be its projection onto the relative class of $\varphi$.
Then $Q\in\mathcal E_\varphi$, $Q\leq z$, and
\[
 \mu_Q\leq\mathbf1_{\{Q=z\}}\mu_z.
\]
\end{proposition}
\begin{proof}
For $C\geq0$ put $p_C=P_\theta(\varphi+C,z)$. The preceding proposition gives $p_0\in\mathcal E_\varphi$. Since $p_0\leq p_C\preceq\varphi$, upward stability gives $p_C\in\mathcal E_\varphi$. The envelope identity, together with $p_C\leq z$, yields the uniform bound
\[
 p_C\leq\varphi+\sup_Xp_C\leq\varphi+\sup_Xz.
\]
Consequently the increasing regularized limit $Q$ is not identically $-\infty$, belongs to $\mathcal E_\varphi$ and is at most $z$. The monotone convergence holds in capacity by \cref{lem:monotone-psh-capacity}.

Finally, \cref{lem:relative-obstacle-contact}, applied with $u=\varphi$ and $w=z$, gives
\begin{equation}\label{eq:big-projection-contact-tail}
 \mu_Q\leq\mathbf1_{\{Q=z\}}\mu_z.
\end{equation}
\end{proof}

\Needspace{16\baselineskip}
\begin{corollary}\label{cor:big-relative-supersolution}
For every nonzero finite non-pluripolar measure $\nu$, there is $Q\in\mathcal E_\varphi$ such that
\[
 0\ne \mathrm{e}^{-Q}\mu_Q\leq\nu.
\]
\end{corollary}
\begin{proof}
By \cref{thm:big-minimal-general-measures}, there is $z\in\mathcal E_V$ with $\mu_z=\mathrm{e}^z\nu$. Since $V$ is an ordinary model and $V\geq\varphi$, the preceding proposition applies with $\eta=V$ and gives $Q=P_\theta[\varphi](z)\in\mathcal E_\varphi$. On the contact set $\{Q=z\}$,
\[
 \mu_Q\leq\mathbf1_{\{Q=z\}}\mathrm{e}^z\nu\leq \mathrm{e}^Q\nu.
\]
The weights are finite off the pole sets, which carry no mass. Thus $\mathrm{e}^{-Q}\mu_Q\leq\nu$, and this measure is nonzero because $\mu_Q(X)\geq\ell>0$.
\end{proof}

\subsection{Varying model references}
\label{subsec:big-varying-models}

We now let the reference vary. Recall from \cref{def:model-volume-comparison} the comparison
\[
 d_\theta(u,v)=2\,\overline{\operatorname{vol}}_\theta(\max(u,v))
 -\overline{\operatorname{vol}}_\theta(u)-\overline{\operatorname{vol}}_\theta(v)
\]
and, for $\delta>0$, the set $\mathcal M_\delta(\theta)$ of ordinary models $u=P_\theta[u]$ with $\underline{\operatorname{vol}}_\theta(u)\geq\delta$. For models $\varphi_j\to\varphi$ in this comparison, \cref{thm:model-volume-capacity-convergence} provides increasing lower models whose upper volumes approach that of $\varphi$. Projecting finite families of potentials onto these lower models makes their rooftops comparable. To keep the limit in the relative class of $\varphi$, domination of the measures is needed only where the potentials lie below one fixed member of that class; domination on all of $X$ gives, in addition, convergence in capacity and convergence of the measures.

\begin{theorem}\label{thm:big-varying-model-closure}\index{Relative full-mass class!varying models}
Let $\varphi_j,\varphi\in\mathcal M_\delta(\theta)$ for some $\delta>0$, with $d_\theta(\varphi_j,\varphi)\to0$. Suppose $u_j\in\mathcal E_{\varphi_j}$, $u_j\to u\in\PSH(X,\theta)$ in $L^1(X)$, and $\chi\in\mathcal E_\varphi$. If one finite non-pluripolar measure $\nu$ satisfies
\[
 \mathbf1_{\{u_j\leq\chi\}}\mu_{u_j}\leq\nu,
\]
then $u\in\mathcal E_\varphi$. Moreover, after extraction and relabeling, the potentials
\[
 \psi_j=P_\theta\left(\inf_{l\geq j}\varphi_l\right),\qquad
 h_j=P_\theta\left(\inf_{l\geq j}u_l\right)
\]
satisfy $\psi_j\in\mathcal M_{\delta/2}(\theta)$ and $h_j\in\mathcal E_{\psi_j}$.
If $\mu_{u_j}\leq\nu$ on $X$, the original sequence converges to $u$ in capacity and $\mu_{u_j}\rightharpoonup\mu_u$.
\end{theorem}
\begin{proof}
We first construct the lower models $\psi_j$, then a common lower bound for rooftops of the $u_l$ in the classes $\mathcal E_{\psi_j}$, and finally treat domination on all of $X$.

By Hartogs' lemma, $u_j$, $u$ and $\chi$ have a common upper bound; translate them to be nonpositive. The envelope identity then gives $u_j\leq\varphi_j$. Since $\varphi_j\to\varphi$ in $L^1$ by \cref{prop:model-volume-separation}, it follows that $u\leq\varphi$. Put $\Lambda=\nu+\mu_\chi$. Using \cref{lem:fixed-npp-integrals}, extract a subsequence along which $u_j\to u$ almost everywhere for smooth volume and
\[
 I_j:=\int_X|\mathrm{e}^{u_j}-\mathrm{e}^u|\,\mathrm{d}\Lambda\leq2^{-j}.
\]
Apply \cref{thm:model-volume-capacity-convergence} to extract further, keeping these estimates, and obtain increasing models $w_j\leq\varphi_j,\varphi$ with lower volume at least $\delta/2$ and $d_\theta(w_j,\varphi)\to0$. For $l\geq j$ we have $w_j\leq w_l\leq\varphi_l$. Hence $\psi_j$ is not identically $-\infty$, and
\[
 w_j\leq\psi_j\leq\varphi.
\]
The upper bound follows by passing to the $L^1$ limit in $\psi_j\leq\varphi_l$ as $l\to\infty$. Monotonicity of the ordinary envelope gives $P_\theta[\psi_j]\leq\varphi_l$ for all $l\geq j$, so $\psi_j$ is an ordinary model. It increases with $j$, has lower volume at least $\delta/2$, and
\begin{equation}\label{eq:varying-reference-volume-limit}
 \overline{\operatorname{vol}}_\theta(\psi_j)
 \longrightarrow\overline{\operatorname{vol}}_\theta(\varphi).
\end{equation}

Set $g_j=P_\theta[\psi_j](\chi)$. \Cref{prop:big-relative-projection}, with $(\varphi,\eta)$ replaced by $(\psi_j,\varphi)$, gives $g_j\in\mathcal E_{\psi_j}$, $g_j\leq\chi$ and $\mu_{g_j}\leq\mu_\chi$. Moreover $g_j$ increases with $j$. For $N\geq j$, consider
\[
 q_{j,N}=P_\theta(g_j,u_j,\ldots,u_N).
\]
For each $l\geq j$, the projection $P_\theta[\psi_j](u_l)$ belongs to $\mathcal E_{\psi_j}$ by the same proposition, now with upper reference $\varphi_l$. The finite rooftop of these projections and $g_j$ lies below $q_{j,N}$. Since $q_{j,N}\leq g_j$, the finite-rooftop proposition and upward stability give $q_{j,N}\in\mathcal E_{\psi_j}$, and hence $\mu_{q_{j,N}}(X)\geq\delta/2$.

Partition the contact set of $q_{j,N}$ into disjoint pieces, one for each obstacle. On a piece where $q_{j,N}=u_l$, we have $u_l\leq g_j\leq\chi$, so the hypothesis bounds the measure there by $\nu$. On the piece where $q_{j,N}=g_j$, the measure is bounded by $\mu_{g_j}\leq\mu_\chi$. Consequently $\mu_{q_{j,N}}\leq\Lambda$. Put $F_j=\min(\mathrm{e}^u,\mathrm{e}^{g_j})$. On the $u_l$ piece,
\[
 |\mathrm{e}^{q_{j,N}}-F_j|\leq|\mathrm{e}^{u_l}-\mathrm{e}^u|.
\]
On the $g_j$ piece, $g_j=q_{j,N}\leq u_N$, so the error is
$(\mathrm{e}^{g_j}-\mathrm{e}^u)_+\leq|\mathrm{e}^{u_N}-\mathrm{e}^u|$.
The $u_N$ piece and the $g_j$ piece are controlled by the two summands $\nu$ and $\mu_\chi$ of $\Lambda$, so together they contribute at most $I_N$. Thus
\begin{equation}\label{eq:varying-roof-exponential-error}
 \int_X|\mathrm{e}^{q_{j,N}}-F_j|\,\mu_{q_{j,N}}
 \leq\sum_{l=j}^N I_l\leq2^{1-j}.
\end{equation}

Since $F_j\geq F_1>0$ outside a pluripolar set, we may choose $\epsilon>0$ with $\Lambda(\{F_1\leq\epsilon\})<\delta/4$. The lower mass bound and the domination by $\Lambda$ give
\[
 \epsilon\delta/4
 \leq\int_X F_j\,\mu_{q_{j,N}}
 \leq2^{1-j}+\Lambda(X)\mathrm{e}^{\sup_Xq_{j,N}}.
\]
For every sufficiently large $j$, this bounds $\sup_Xq_{j,N}$ from below, uniformly in $N$. Hence the decreasing limit $q_j$ as $N\to\infty$ is not identically $-\infty$. Applying \cref{thm:big-relative-fixed-continuity} in the fixed class $\mathcal E_{\psi_j}$, with dominating measure $\Lambda$, gives $q_j\in\mathcal E_{\psi_j}$. Now
\[
 q_j\leq h_j\leq\psi_j,
\]
so $h_j\in\mathcal E_{\psi_j}$ by upward stability. We discard the finitely many earlier indices. Since $h_j\leq u\leq\varphi$, monotonicity and \eqref{eq:varying-reference-volume-limit} show that $u$ and $\varphi$ have the same upper volume, and \cref{thm:relative-equivalence} gives $u\in\mathcal E_\varphi$.

Suppose now that $\mu_{u_j}\leq\nu$ on $X$. The finite rooftops $p_{j,N}=P_\theta(u_j,\ldots,u_N)$ are not identically $-\infty$, since they lie above $h_j$. Partitioning their contact sets gives
\[
 \mu_{p_{j,N}}\leq\nu,\qquad
 \int_X|\mathrm{e}^{p_{j,N}}-\mathrm{e}^u|\,\mu_{p_{j,N}}\leq2^{1-j}.
\]
For fixed $j$, they decrease to $h_j$ as $N\to\infty$. These lower envelopes increase to a regularized limit $h\leq u\leq\varphi$, and their upper volumes tend to that of $\varphi$; hence $h\in\mathcal E_\varphi$. Both monotone limits hold in capacity by \cref{lem:monotone-psh-capacity}. Applying \cref{lem:weighted-extended-lsc} twice gives
\[
 \int_X|\mathrm{e}^h-\mathrm{e}^u|\,\mu_h=0.
\]
Hence $\mu_h(\{h<u\})=0$, and zero domination gives $h=u$. The upper envelopes $(\sup_{k\geq j}u_k)^*$ decrease to $u$ by Hartogs' lemma and the almost-everywhere convergence, and squeezing gives convergence in capacity. Finally,
\[
 \mu_{u_j}(\{u_j\leq-T\})\leq\nu(\{h_1\leq-T\})\longrightarrow0
\]
uniformly in $j$, and the tail criterion \cref{thm:general-background-damping} gives weak convergence of the measures. Since the extraction can start from any subsequence, both conclusions hold for the original sequence.
\end{proof}


\section{General measures with a model reference}
\label{sec:big-model-general-measures}

The projection produces a nonzero weighted measure below the prescribed datum. Maximize its mass among supersolutions with that same upper bound. If any part of the datum remains, weighted addition increases the mass and contradicts maximality. Density truncation then gives the equation for every finite non-pluripolar measure.

\begin{theorem}\label{thm:big-model-exponential}
Let $\mu$ be a nonzero finite non-pluripolar measure and let $\lambda>0$. There is a unique $u\in\mathcal E_\varphi$ such that
\[
 \mathrm{MA}_\theta(u)=\mathrm{e}^{\lambda u}\mu.
\]
\end{theorem}
\begin{proof}
We first take $\lambda=1$ and suppose that
\[
 \mu\leq A\sigma,\qquad A>0,\qquad
 \sigma=\mathrm{MA}_\omega(\psi),\quad
 \psi\in\PSH(X,\omega)\cap L^\infty(X).
\] The supersolution corollary gives $Q\in\mathcal E_\varphi$ with a finite nonzero measure
\[
 \nu_0=\mathrm{e}^{-Q}\mu_Q\leq\mu.
\]
Consider
\[
 \mathcal T=\{w\in\mathcal E_\varphi:
              \nu_0\leq \mathrm{e}^{-w}\mu_w\leq\mu\}.
\]
The family contains $Q$. Exponential order gives $w\leq Q$ for every $w\in\mathcal T$. Also
\[
 \ell\leq\mu_w(X)\leq \mathrm{e}^{\sup_Xw}\mu(X),
 \qquad
 \sup_Xw\geq\log\frac{\ell}{\mu(X)}.
\]
These bounds prevent collapse of any sequence in $\mathcal T$.

\Needspace{6\baselineskip}
Let
\[
 \alpha=\sup_{w\in\mathcal T}\int_X\mathrm{e}^{-w}\,\mu_w\leq\mu(X),
\]
and choose $w_j\in\mathcal T$ whose weighted masses tend to $\alpha$. Extract a genuine $L^1$ limit $w$. Since
\[
 \mu_{w_j}\leq \mathrm{e}^{\sup_XQ}\mu,
\]
the fixed-dominator theorem gives $w\in\mathcal E_\varphi$, convergence in capacity, and $\mu_{w_j}\rightharpoonup\mu_w$. Put $\nu_j=\mathrm{e}^{-w_j}\mu_{w_j}$. After another extraction, these finite measures converge weakly to a measure $\nu$, and
\[
 \nu_0\leq\nu\leq\mu,\qquad \nu(X)=\alpha.
\]

We claim $\mu_w=\mathrm{e}^w\nu$. \Cref{lem:fixed-npp-integrals} gives
\[
 \int_X|\mathrm{e}^{w_j}-\mathrm{e}^w|\,\mathrm{d}\mu\longrightarrow0.
\]
If $f$ is continuous, replacing $\mathrm{e}^{w_j}$ by $\mathrm{e}^w$ in $\int_Xf \mathrm{e}^{w_j}\,\mathrm{d}\nu_j$ therefore has error tending to zero. For the remaining integral, approximate the fixed bounded measurable function $f \mathrm{e}^w$ in $L^1(\mu)$ by continuous functions. Since $\nu_j,\nu\leq\mu$, the approximation error is uniform in $j$, while the integrals of each continuous approximation converge by weak convergence. Thus
\[
 \int_Xf\,\mu_{w_j}
 =\int_Xf \mathrm{e}^{w_j}\,\mathrm{d}\nu_j
 \longrightarrow\int_Xf \mathrm{e}^w\,\mathrm{d}\nu.
\]
This proves the claim. The measure $\nu$ ignores the poles of $w$, so $\mathrm{e}^{-w}\mu_w=\nu$. Hence $w\in\mathcal T$ attains the maximum $\alpha$.

If $\mu-\nu$ were nonzero, the supersolution corollary would supply $v\in\mathcal E_\varphi$ with
\[
 0\ne\nu_v=\mathrm{e}^{-v}\mu_v\leq\mu-\nu.
\]
In particular $\mu_v\leq \mathrm{e}^{\sup_Xv}\mu\leq A \mathrm{e}^{\sup_Xv}\sigma$, so the bounded-carrier hypothesis of the addition theorem is satisfied. That theorem gives $a\in\mathcal E_\varphi$ with
\[
 \mathrm{e}^{-a}\mu_a=\mathrm{e}^{-w}\mu_w+\mathrm{e}^{-v}\mu_v=\nu+\nu_v.
\]
The right side is between $\nu_0$ and $\mu$, but its mass exceeds $\alpha$. This contradicts maximality. Therefore $\nu=\mu$, and $w$ solves the equation.

\Needspace{8\baselineskip}
Now let $\mu$ be any nonzero finite non-pluripolar measure. The bounded-carrier decomposition writes
\[
 \mu=f\sigma,\qquad
 \sigma=\mathrm{MA}_\omega(\psi),\quad
 \psi\in\PSH(X,\omega)\cap L^\infty(X),\quad
 0\leq f\in L^1(\sigma).
\]
\Needspace{8\baselineskip}
Set $f_k=\min(f,k)$ and $\mu_k=f_k\sigma$. These measures are nonzero, and the first part gives $u_k\in\mathcal E_\varphi$ with $\mu_{u_k}=\mathrm{e}^{u_k}\mu_k$. Exponential order shows that $u_k$ decreases. The uniform bound
\[
 \sup_Xu_k\geq\log\frac{\ell}{\mu(X)}
\]
gives a genuine limit $u$. Moreover $\mu_{u_k}\leq \mathrm{e}^{\sup_Xu_1}\mu$. Fixed-dominator continuity gives $u\in\mathcal E_\varphi$ and weak convergence of these measures. On the right,
$\mathrm{e}^{u_k}f_k\to \mathrm{e}^u f$ in $L^1(\sigma)$ by dominated convergence, with dominator $\mathrm{e}^{\sup_Xu_1}f$. Thus $\mu_u=\mathrm{e}^u\mu$.

For general $\lambda>0$, apply the exponent-one result in background $\lambda\theta$, with model $\lambda\varphi$ and datum $\lambda^n\mu$, and divide its potential by $\lambda$. Scaling of the class and the Monge--Amp\`ere measure gives the displayed equation. Finally, if $u,v\in\mathcal E_\varphi$ solve it, then $\mu_u=\mathrm{e}^{\lambda(u-v)}\mu_v$. The strict domination theorem with constant shifts gives $u=v$.
\end{proof}

\Needspace{15\baselineskip}
\begin{theorem}\label{thm:big-model-normalized}
For every nonzero finite non-pluripolar measure $\mu$, there are $u\in\mathcal E_\varphi$ and $c>0$ such that
\[
 \sup_Xu=0,\qquad \mathrm{MA}_\theta(u)=c\mu.
\]
The scalar $c$ is unique and satisfies
\[
 \frac{\ell}{\mu(X)}\leq c\leq\frac{M}{\mu(X)}.
\]
\end{theorem}
\begin{proof}
For each integer $j\geq1$, solve
\[
 \mu_{u_j}=\mathrm{e}^{u_j/j}\mu,\qquad u_j\in\mathcal E_\varphi,
\]
using the preceding theorem. Put $S=\sup_Xu_1$ and $w_j=u_1+(j-1)S$. Then $w_j/j\geq u_1$, and hence
\[
 \mu_{w_j}=\mathrm{e}^{u_1}\mu\leq \mathrm{e}^{w_j/j}\mu.
\]
Exponential order at exponent $1/j$ gives $u_j\leq w_j$. Define
\[
 v_j=u_j-\sup_Xu_j,\qquad
 c_j=\mathrm{e}^{(\sup_Xu_j)/j}.
\]
The upper bound just proved and the positive lower mass imply
\[
 \frac{\ell}{\mu(X)}\leq c_j\leq \mathrm{e}^S,
 \qquad \mu_{v_j}=c_j \mathrm{e}^{v_j/j}\mu\leq \mathrm{e}^S\mu.
\]
The envelope identity gives $v_j\leq\varphi$. \Cref{thm:normalized-compactness} supplies a subsequence $j_k$ for which $v_{j_k}\to u$ in $L^1$, with $u$ genuine and $\sup_Xu=0$, and $c_{j_k}\to c>0$. Fixed-dominator continuity gives $u\in\mathcal E_\varphi$ and $\mu_{v_{j_k}}\rightharpoonup\mu_u$.

\Cref{lem:normalized-exponential-npp} gives $\mathrm{e}^{v_j/j}\to1$ in $L^1(\mu)$, also for this non-semipositive background. Consequently $c_{j_k}\mathrm{e}^{v_{j_k}/j_k}\mu\to c\mu$ in total variation. The retained reciprocal exponents are $1/j_k$, so the limit equation is $\mu_u=c\mu$. Integrating it gives the stated interval for $c$. The scalar comparison in \cref{cor:big-relative-orders} gives uniqueness of $c$.
\end{proof}

For $L^p$ data, \cref{sec:big-model-lp} will prove the stronger conclusion $u\simeq\varphi$.

\subsection{Stability under changes of reference}
\label{subsec:big-equation-reference-stability}

Convergence of the models in $d_\theta$ controls the relative class of a limit. For equations, convergence of the data supplies the fixed dominating measures needed after extraction. The positive exponent then identifies the potential; at exponent zero, scalar uniqueness identifies the limiting measure.

\begin{theorem}\label{thm:big-varying-equation-stability}
Let $\varphi_j,\varphi\in\mathcal M_\delta(\theta)$ for some $\delta>0$, with $d_\theta(\varphi_j,\varphi)\to0$. Let $\nu$ be a finite non-pluripolar measure and let $f_j,f$ be nonnegative nonzero functions in $L^1(\nu)$ with $f_j\to f$ in $L^1(\nu)$. Put $\mu_j=f_j\nu$ and $\mu=f\nu$.
\begin{enumerate}
\item For a fixed $\lambda>0$, the solutions
\[
 u_j\in\mathcal E_{\varphi_j},\qquad
 \mathrm{MA}_\theta(u_j)=\mathrm{e}^{\lambda u_j}\mu_j
\]
converge in $L^1(X)$ and in capacity to the solution $u\in\mathcal E_\varphi$ with datum $\mu$. Their Monge--Amp\`ere measures converge to $\mathrm{MA}_\theta(u)$ in total variation.
\item Let $u_j\in\mathcal E_{\varphi_j}$ satisfy $\sup_Xu_j=0$ and $\mathrm{MA}_\theta(u_j)=c_j\mu_j$. Then $c_j\to c$, where $c$ is the scalar for reference $\varphi$ and datum $\mu$, and
\[
 \mathrm{MA}_\theta(u_j)\longrightarrow c\mu
 \quad\text{in total variation}.
\]
The potentials form an $L^1$ relatively compact family. Every $L^1$ cluster limit $u$ lies in $\mathcal E_\varphi$, has supremum zero and solves $\mathrm{MA}_\theta(u)=c\mu$. Each subsequence converging in $L^1$ converges to that limit in capacity.
\end{enumerate}
\end{theorem}
\begin{proof}
Write $M_0=\overline{\operatorname{vol}}_\theta(V_\theta)$. Every solution above has mass between $\delta$ and $M_0$. In the first case let $s_j=\sup_Xu_j$. Then
\[
 s_j\geq\lambda^{-1}\log\frac{\delta}{\mu_j(X)},
\]
which is a uniform lower bound. To bound $s_j$ above, suppose after extraction that $s_j\to+\infty$. Normalized compactness gives $v_j=u_j-s_j\to v$ in $L^1$ with $\sup_Xv=0$. The fixed-measure lemma and $f_j\to f$ in $L^1(\nu)$ imply
\[
 \int_X\mathrm{e}^{\lambda v_j}f_j\,\mathrm{d}\nu
 \longrightarrow\int_X\mathrm{e}^{\lambda v}f\,\mathrm{d}\nu>0.
\]
The limit is positive since $f\nu$ is nonzero and ignores the poles of $v$. But the left side equals $\mathrm{e}^{-\lambda s_j}\mu_{u_j}(X)\leq M_0\mathrm{e}^{-\lambda s_j}$, a contradiction. Thus $s_j$ is bounded and the $u_j$ form an $L^1$ relatively compact family of genuine potentials.

Start with any subsequence and extract $u_j\to u$ in $L^1$. Extract again so that $\sum_j\|f_j-f\|_{L^1(\nu)}<\infty$. The function
\[
 g=f+\sum_j|f_j-f|\in L^1(\nu)
\]
dominates all retained $f_j$. If $C$ is their common upper bound for $u_j$, then $\mu_{u_j}\leq\mathrm{e}^{\lambda C}g\nu$. By \cref{thm:big-varying-model-closure}, $u\in\mathcal E_\varphi$, the potentials converge in capacity and the measures converge weakly to $\mu_u$. On the other hand,
\[
 \mathrm{e}^{\lambda u_j}f_j\longrightarrow
 \mathrm{e}^{\lambda u}f\quad\text{in }L^1(\nu).
\]
Indeed, the error from $f_j-f$ is at most $\mathrm{e}^{\lambda C}\|f_j-f\|_1$, and \cref{lem:fixed-npp-integrals} for the fixed measure $f\nu$ controls the other term. Hence $\mu_u=\mathrm{e}^{\lambda u}\mu$, with convergence in total variation. Uniqueness in \cref{thm:big-model-exponential} identifies every cluster limit. The subsequence criterion gives all three convergences for the original sequence.

For the second case, integrating the equation gives
\[
 \frac{\delta}{\mu_j(X)}\leq c_j\leq\frac{M_0}{\mu_j(X)}.
\]
Since $\mu_j(X)\to\mu(X)>0$, these scalars stay in a compact subinterval of $(0,+\infty)$. Normalized compactness gives, after extraction, $u_j\to u$ in $L^1$, $\sup_Xu=0$, and $c_j\to c'>0$. Extract a summable density error as above. Then $\mu_{u_j}\leq Cg\nu$ for a fixed $C$, so \cref{thm:big-varying-model-closure} gives $u\in\mathcal E_\varphi$, capacity convergence and weak convergence to $\mu_u$. Since $c_jf_j\to c'f$ in $L^1(\nu)$, the limit equation is $\mu_u=c'\mu$. Scalar uniqueness in \cref{thm:big-model-normalized} gives $c'=c$. Thus the whole scalar sequence converges, and so does $c_j\mu_j$ in total variation. Applying the same argument within any $L^1$ convergent subsequence proves its capacity convergence and the asserted cluster equation.
\end{proof}


\subsection{Arbitrary weighted sums}
\label{sec:big-unrestricted-addition}

With existence for general measures available, weighted addition extends to arbitrary $u,v\in\mathcal E_\varphi$. Truncating the exponential weights gives finite data, which \cref{thm:big-model-exponential} solves. A first-moment estimate prevents collapse of these solutions, and a rooftop involving their limit localizes the equation before the limit can be compared with the inputs.

\Needspace{11\baselineskip}
\begin{theorem}\label{thm:big-unrestricted-addition}\index{Monge--Ampere measure@Monge--Amp\`ere measure!addition in relative classes}
For every $u,v\in\mathcal E_\varphi$ and $\lambda>0$, there is a unique $a\in\mathcal E_\varphi$ satisfying
\[
 \mu_a=\mathrm{e}^{\lambda(a-u)}\mu_u+\mathrm{e}^{\lambda(a-v)}\mu_v.
\]
Moreover $a\leq\min(u,v)$ and $\mu_a\leq\mu_u+\mu_v$.
\end{theorem}
\begin{proof}
By the scaling and the common translation used in \cref{thm:big-weighted-addition}, it suffices to take $\lambda=1$ and $u,v\leq0$. For positive integers $j$, let
\[
 \tau_j=\mathrm{e}^{-\max(u,-j)}\mu_u+\mathrm{e}^{-\max(v,-j)}\mu_v.
\]
These are nonzero finite non-pluripolar measures, and they increase with $j$. Only the weights are truncated; the functions $\max(u,-j)$ and $\max(v,-j)$ need not be $\theta$-plurisubharmonic. \Cref{thm:big-model-exponential} gives $a_j\in\mathcal E_\varphi$ with $\mu_{a_j}=\mathrm{e}^{a_j}\tau_j$. Since $\tau_j$ increases, exponential order implies that $a_j$ decreases.

Set $r=P_\theta(u,v)\in\mathcal E_\varphi$. On $\{a_j<r\}$, one has $r\leq u,v$ and
\[
 (r-a_j)\mathrm{e}^{a_j-\max(u,-j)}\leq(r-a_j)\mathrm{e}^{a_j-r}\leq \mathrm{e}^{-1},
\]
with the same inequality for $v$. Consequently
\[
 \int_X(r-a_j)_+\,\mu_{a_j}
 \leq \mathrm{e}^{-1}\bigl(\mu_u(X)+\mu_v(X)\bigr).
\]
Also $a_j\leq a_1\leq\varphi+\sup_Xa_1$. The one-sided first-moment criterion of \cref{prop:big-relative-first-moment}, with fixed reference $r$, bounds $\sup_Xa_j$ from below. Thus the decreasing sequence has an $L^1$ limit $a\not\equiv-\infty$, and the same criterion puts $a$ in $\mathcal E_\varphi$.

We identify the equation before comparing $a$ with the inputs. The rooftop
\[
 q=P_\theta(a,u,v)\in\mathcal E_\varphi
\]
lies below $a_j$, $a$, $u$ and $v$.

\Needspace{10\baselineskip}
For $T>0$, let $\chi_T=g_T(q)$, with $g_T$ as in the proof of \cref{lem:big-truncated-addition-limit}: a continuous nondecreasing cutoff that vanishes where $q\leq-2T$, equals one where $q\geq-T$, and increases to one off the pole set of $q$ as $T\to\infty$. On $\{\chi_T>0\}$ both inputs exceed $-2T$. Hence, for $j>2T$,
\begin{equation}\label{eq:unrestricted-addition-cutoff}
 \chi_T\mu_{a_j}
 =\chi_T \mathrm{e}^{a_j}(\mathrm{e}^{-u}\mu_u+\mathrm{e}^{-v}\mu_v).
\end{equation}
The factors $\chi_T\mathrm{e}^{-u}$ and $\chi_T\mathrm{e}^{-v}$ are bounded by $\mathrm{e}^{2T}$. Since $a_j\downarrow a$ and $a_j\leq\sup_Xa_1$, dominated convergence against the fixed measure $\mu_u+\mu_v$ gives convergence in total variation of the right-hand sides to the corresponding expression with $a$.

For the left-hand side, $a_j\to a$ in capacity by \cref{lem:monotone-psh-capacity}. Choose a Hermitian form $\Omega\geq\theta$, and replace $a_j$ by $b_j=\max(a_j,-2T)$ before forming bounded products. The $b_j$ are uniformly bounded $\Omega$-plurisubharmonic functions converging in capacity. The expansion
\[
 (\theta+\ddc b_j)^n
 =\sum_{k=0}^n\binom nk(\Omega+\ddc b_j)^k
                         \wedge(\theta-\Omega)^{n-k}
\]
and bounded capacity continuity apply against the fixed quasi-continuous weight $\chi_T$. Where $\chi_T\ne0$ no potential is truncated, so locality gives convergence to $\chi_T\mu_a$. In this way the possibly signed bounded product is handled before the cutoff is inserted.

Thus \eqref{eq:unrestricted-addition-cutoff} holds in the limit. Let $T\to\infty$. Both sides ignore the poles of $q$, so monotone convergence gives
\[
 \mu_a=\mathrm{e}^{a-u}\mu_u+\mathrm{e}^{a-v}\mu_v.
\]
Since $a\in\mathcal E_\varphi$, this also shows that the weighted right-hand side is finite.

For $\varepsilon>0$, the equation implies
\[
 \mathbf1_{\{u<a-\varepsilon\}}\mu_u
 \leq \mathrm{e}^{-\varepsilon}\mathbf1_{\{u<a-\varepsilon\}}\mu_a.
\]
Strict domination, \cref{thm:big-relative-domination}, gives $u\geq a-\varepsilon$, and letting $\varepsilon\downarrow0$ gives $u\geq a$. Similarly $v\geq a$, so $\mu_a\leq\mu_u+\mu_v$. Two solutions have the same exponentially weighted measure, so \cref{cor:big-relative-orders}, applied in both directions, gives uniqueness.
\end{proof}

\begin{corollary}\label{cor:weighted-addition-algebra}
For fixed $\lambda>0$, write $u\star_\lambda v$ for the solution $a$ in the theorem. This operation on $\mathcal E_\varphi$ is commutative and associative, and, for $c\in\mathbb R$,
\[
 (u+c)\star_\lambda(v+c)=(u\star_\lambda v)+c,
 \qquad
 u\star_\lambda u=u-\lambda^{-1}\log2.
\]
\end{corollary}
\begin{proof}
Put $F_\lambda(u)=\mathrm{e}^{-\lambda u}\mu_u$, with zero mass on the poles. The theorem says that
$F_\lambda(u\star_\lambda v)=F_\lambda(u)+F_\lambda(v)$.
By \cref{cor:big-relative-orders}, the map $F_\lambda$ is injective on $\mathcal E_\varphi$. Commutativity and associativity therefore follow from those of addition of positive measures, including measures of infinite total mass. Finally,
$F_\lambda(u+c)=\mathrm{e}^{-\lambda c}F_\lambda(u)$, which together with injectivity gives both identities.
\end{proof}


\Needspace{12\baselineskip}
\section{Relative boundedness for \texorpdfstring{$L^p$}{Lp} data}
\label{sec:big-model-lp}

The general-measure theorem gives a solution in $\mathcal E_\varphi$. We now show that $L^p$ data, with $p>1$, force ordinary singularity type $\varphi$. We keep the same bounded-mass and positive lower reference-volume assumptions. The estimate will be applied directly to a relative-class potential: truncating it by $\max(u,\varphi-j)$ would introduce a contact measure not controlled by the original density.

\begin{lemma}\label{lem:big-model-concave-envelope}
Let $u\in\mathcal E_\varphi$, $u\leq\varphi$, and let
$\chi:\mathbb R_{\leq0}\to\mathbb R_{\leq0}$ be increasing and concave, with $\chi(0)=0$. Suppose it is $C^{1,1}$ on compact intervals and
\[
 1\leq\chi'\leq L<\infty.
\]
Then the envelope
\[
 p=P_\theta(\varphi+\chi(u-\varphi))
\]
belongs to $\mathcal E_\varphi$, and
\[
 \mu_p\leq\mathbf1_{\{p=\varphi+\chi(u-\varphi)\}}
                 \chi'(u-\varphi)^n\mu_u.
\]
The obstacle and the contact set are defined off the polar indeterminacy sets, which all these measures ignore.
\end{lemma}
\begin{proof}
For $t\leq0$ the derivative bounds give $Lt\leq\chi(t)\leq t$. Thus
\[
 Q_L(u,\varphi)\leq p\leq u\leq\varphi.
\]
The lower competitor belongs to $\mathcal E_\varphi$ by stretching, so $p$ is genuine and belongs to that class by upward stability. The obstacle is finite and quasi-continuous off poles, and the contact theorem applies.

Let $\tau=\chi^{-1}$. It is convex, with $L^{-1}\leq\tau'\leq1$. We justify the transformation
\[
 w=\varphi+\tau(p-\varphi)
\]
without assuming any open set on which $\varphi$ is locally bounded. On a coordinate ball choose a smooth $q$ with $\theta+\ddc q\geq0$. The functions
\[
 p_m=\max(p,q-m),\qquad
 \varphi_m=\max(\varphi,q-m)
\]
are locally bounded $\theta$-plurisubharmonic functions, and $p_m\leq\varphi_m$. The bounded convex chain calculation in the proof of \cref{lem:big-concave-envelope} applies to
$w_m=\varphi_m+\tau(p_m-\varphi_m)$ and gives
\[
 \theta_{w_m}\geq\tau'(p_m-\varphi_m)\theta_{p_m}
       +(1-\tau'(p_m-\varphi_m))\theta_{\varphi_m}\geq0,
 \qquad
 \mu_{w_m}\geq\tau'(p_m-\varphi_m)^n\mu_{p_m}.
\]
In that calculation the two current coefficients add to one; no positivity of $\theta$ itself is used. Its smoothing argument applies on each smaller ball to these locally bounded potentials and the continuous derivative factors.

The scalar function $(s,t)\mapsto t+\tau(s-t)$ on $s\leq t$ is nondecreasing in each variable. Consequently $w_m$ decreases as $m$ increases. Since $p\leq p_m\leq w_m\leq\varphi_m$, it has a genuine local limit. Off the poles this limit is exactly $w$ as displayed above. The local limits therefore agree canonically on overlaps and define a global $\theta$-plurisubharmonic function.

On $D_m=\{p>q-m\}$ neither $p$ nor $\varphi$ is truncated. Thus $w_m=w$, and plurifine locality transfers the product inequality to $D_m$. These sets exhaust the complement of the poles of $p$, locally on the ball. All products ignore the discarded polar set, so
\[
 \mu_w\geq\tau'(p-\varphi)^n\mu_p.
\]
The obstacle bound gives $w\leq u$, with equality on the contact set of $p$. Ordered contact, the identity
$\tau'(p-\varphi)=1/\chi'(u-\varphi)$ there, and contact support of $\mu_p$ prove the assertion.
\end{proof}

\begin{proposition}\label{prop:big-model-moment-bound}
Let $\mu$ be a finite non-pluripolar measure. Suppose, for some real $m>2n+1$,
\[
 A_m:=\sup_{\substack{v\in\PSH(X,\theta)\\\sup_Xv=0}}
             \int_X(-v)^m\,\mathrm{d}\mu<\infty.
\]
There is a finite constant $C$, depending only on
$n,\ell,M,m,\mu(X)$ and $A_m$, such that every $u\in\mathcal E_\varphi$ with $\sup_Xu=0$ and $\mu_u\leq\mu$ satisfies
\[
 \varphi-C\leq u\leq\varphi.
\]
\end{proposition}
\begin{proof}
The envelope identity gives $u\leq\varphi\leq0$. Put $x=\varphi-u$, defined off poles, and $N=2n+1$. H\"older's inequality gives the uniform first-moment bound
\[
 B:=\int_Xx\,\mu_u
 \leq\int_X(-u)\,\mathrm{d}\mu
 \leq\mu(X)^{1-1/m}A_m^{1/m}=:A_1.
\]
Let $H(t)=\mu_u(\{x>t\})$ and
$T=\sup\{t\geq0:H(t)>0\}$, with $T=0$ if the set is empty. At this point $T$ may be infinite. If $T<\infty$, zero domination applied to $\varphi-t$ for $t>T$ gives $u\geq\varphi-T$. It is therefore enough to bound $T$. We may assume $T>1$ and choose any $S$ with $1<S<T$.

Define, for $t\geq0$,
\[
 g(t)=
 \begin{cases}
  ((1+t)^{3/2}H(t))^{-1},&0\leq t\leq S,\\
  (1+t)^{-3/2},&t>S,
 \end{cases}
 \qquad F(t)=1+\int_0^t g(s)\,\mathrm{d}s,
\]
and set $W(t)=\int_0^t F(s)^{1/N}\,\mathrm{d}s$ and $\chi(-t)=-W(t)$. Since $H(S)>0$, $F$ is bounded on the entire half-line, although its bound may depend on $S$ and $u$. It is locally Lipschitz and nondecreasing. Thus $\chi$ has bounded derivative, is increasing and concave, and meets the preceding lemma's hypotheses. Also $W(t)\geq t$ and $W(t)\leq tF(t)^{1/N}$.

Tonelli's theorem and $H(t)\leq M$ give a bound independent of $S$ and $u$:
\[
 \int_X F(x)\,\mu_u
 =\mu_u(X)+\int_0^\infty g(t)H(t)\,\mathrm{d}t
 \leq3\max(1,M)=:K.
\]
Let $p=P_\theta(\varphi-W(x))\in\mathcal E_\varphi$. On its contact set the preceding lemma bounds its measure by $F(x)^{n/N}\mu_u$. Hence, using $N=2n+1$,
\[
 \begin{split}
 \int_X(\varphi-p)^{1/2}\,\mu_p
 &\leq\int_X x^{1/2}F(x)^{(n+1/2)/N}\,\mu_u\\
 &\leq\left(\int_Xx\,\mu_u\right)^{1/2}
          \left(\int_XF(x)\,\mu_u\right)^{1/2}
 \leq\sqrt{KB}.
 \end{split}
\]
The identity $P_\theta[p]=\varphi$ implies $p-\sup_Xp\leq\varphi$. Since $p\leq\varphi\leq0$, the last integral is at least $\ell\sqrt{-\sup_Xp}$. Consequently
\[
 0\leq-\sup_Xp\leq KA_1/\ell^2=:D.
\]
Using $W(x)\leq\varphi-p\leq-p$, and applying the assumed moment bound to $p-\sup_Xp$, we obtain
\begin{equation}\label{eq:big-model-weight-moment}
 \int_XW(x)^m\,\mathrm{d}\mu
 \leq2^{m-1}\bigl(A_m+D^m\mu(X)\bigr)<C_0,
\end{equation}
where we choose $C_0=1+2^{m-1}(A_m+D^m\mu(X))$. This constant is independent of $S$ and $u$.

For $0<t\leq S$, Chebyshev's inequality and $\mu_u\leq\mu$ give
\[
 H(t)\leq C_0W(t)^{-m},
 \qquad
 F'(t)\geq\frac{W(t)^m}{C_0(1+t)^{3/2}}
 \quad\text{for almost every }t.
\]
Since $W'=F^{1/N}$, the derivative of $(W')^{N+1}$ is at least
\[
 \frac{N+1}{NC_0}\,
 \frac{W^mW'}{(1+t)^{3/2}}.
\]
All functions being differentiated are absolutely continuous on compact intervals. Integrating from zero to $t$, using $W(0)=0$, gives
\[
 (W'(t))^{N+1}
 \geq\frac{N+1}{NC_0(m+1)}
        \frac{W(t)^{m+1}}{(1+t)^{3/2}}.
\]
In particular
\[
 W'(t)\geq c_0(1+t)^{-\gamma}W(t)^\kappa,
 \quad
 \gamma=\frac{3}{2(N+1)}<1,\quad
 \kappa=\frac{m+1}{N+1}>1,
\]
with $c_0>0$ independent of $S,u$. Integrating the derivative of $W^{1-\kappa}$ on $[1,S]$ yields
\[
 (\kappa-1)c_0\int_1^S(1+t)^{-\gamma}\,\mathrm{d}t
 \leq W(1)^{1-\kappa}-W(S)^{1-\kappa}\leq1,
\]
because $W(1)\geq1$. The integral diverges as $S\to\infty$. It follows that every $S<T$ has a common finite upper bound, and hence so does $T$. The zero-domination argument at the start now gives $u\geq\varphi-C$. The cases $T\leq1$, including $B=0$, were already covered.
\end{proof}

\begin{corollary}\label{cor:big-model-lp-type}
Let $\mathrm{d}V$ be a smooth positive volume form and let $0\leq f\in L^p(X,\mathrm{d}V)$, with $p>1$ and $\int_Xf\,\mathrm{d}V>0$.
For $\lambda>0$, the solution in \cref{thm:big-model-exponential} to
\[
 \mathrm{MA}_\theta(u)=e^{\lambda u}f\,\mathrm{d}V
\]
satisfies $u\simeq\varphi$. Every normalized solution in $\mathcal E_\varphi$ of
\[
 \mathrm{MA}_\theta(u)=cf\,\mathrm{d}V,\qquad c>0,
\]
also satisfies $u\simeq\varphi$. In particular both existence statements can be made in the ordinary singularity type of $\varphi$.
\end{corollary}
\begin{proof}
Uniform Skoda integrability for normalized $\theta$-plurisubharmonic functions gives uniform moments of every finite order against $\mathrm{d}V$. H\"older's inequality therefore implies, for every finite $m$,
\[
 \sup_{\substack{v\in\PSH(X,\theta)\\\sup_Xv=0}}
       \int_X(-v)^m f\,\mathrm{d}V<\infty.
\]
Choose $m>2n+1$. For the positive-exponent solution, set $S=\sup_Xu$ and $v=u-S$. Then $v\in\mathcal E_\varphi$, $\sup_Xv=0$, and
\[
 \mu_v=e^{\lambda u}f\,\mathrm{d}V
 \leq e^{\lambda S}f\,\mathrm{d}V.
\]
The proposition applied to this fixed dominating measure shows that $v$, and hence $u$, has ordinary type $\varphi$. For the normalized equation apply it to the measure $cf\,\mathrm{d}V$. This proves the ordinary-type conclusion of \cite[Corollary 6.1]{XL26} directly from the general-measure solutions and the relative estimate above.
\end{proof}

The derivative bound on the concave weight was allowed to depend on the chosen depth $S$. The moment estimate, the normalization of its envelope, and the final depth bound were independent of that derivative bound. This is what permits the argument to start with $u\in\mathcal E_\varphi$ without first assuming $u\simeq\varphi$.


\Needspace{24\baselineskip}
\section{Comparison of the equation regimes}

Let $X$ be a connected compact complex manifold and let $\theta$ be a smooth real $(1,1)$-form on $X$. Write $\mathrm{MA}_\theta(a)$ for its extended non-pluripolar measure when $a\in\PSH(X,\theta)$, and $a\preceq b$ when $a\leq b+C$ for some real constant $C$. The operators $P_\theta$, $P_\theta[\,\cdot\,]$ and the two relative classes are defined in \cref{chap:envelopes,chap:mass-bound-volumes}. A model reference $\varphi\in\PSH(X,\theta)$ is specified only in the regimes that use it.

For comparison, the following statements use the hypotheses listed in each regime. The first regime does not require either uniform mass property; the second requires positive lower mass; the third also requires bounded mass.

\Needspace{18\baselineskip}
\begin{theorem}\label{thm:relative-domination}
Let $\theta$ be a smooth real $(1,1)$-form, let $u,v\in\PSH(X,\theta)$, and let $0\leq a<1$. If
\[
 \mathbf1_{\{u<v\}}\mathrm{MA}_\theta(u)
 \leq a\,\mathbf1_{\{u<v\}}\mathrm{MA}_\theta(v),
\]
then $u\geq v$ in each of the following regimes:
\begin{enumerate}
\item $\theta$ is semipositive and big, $P_{\delta\theta}(u)\not\equiv-\infty$ for some $\delta\in(0,1)$, and $v\preceq u$;
\item $\theta,\varphi$ satisfy the assumptions of \cref{thm:als-general-measure}, $u\in\mathcal E_{\mathrm{env}}(X,\theta,\varphi)$, $v\preceq\varphi$;
\item $X$ has bounded mass, $\theta$ is big, $\varphi=P_\theta[\varphi]$ has positive lower volume, $u\in\mathcal E_{\mathrm{vol}}(X,\theta,\varphi)$, and $v\preceq\varphi$.
\end{enumerate}
\end{theorem}

The first assertion was proved in \cref{thm:singular-strict-domination}. The second is proved without a finite-mass assumption in \cref{thm:relative-extended-domination}. The third is proved in \cref{thm:big-relative-domination}. In the first case, a potential equivalent to an ALS model meets the smaller-background condition by a constant translation of a competitor for that model.

\begin{proposition}\label{prop:relative-uniqueness}
Fix a nonzero finite positive non-pluripolar measure $\mu$. Consider either a fixed ordinary singularity type admitting a smaller-background envelope in a semipositive big background, or one of the relative classes in parts (2) and (3) of \cref{thm:relative-domination}. In each such class, two solutions of $\mathrm{MA}_\theta(w)=\mathrm{e}^{\lambda w}\mu$ with $\lambda>0$ coincide. If $u$ and $v$ solve
\[
 \mathrm{MA}_\theta(u)=c_u\mu,\qquad
 \mathrm{MA}_\theta(v)=c_v\mu,
 \qquad c_u,c_v>0,
\]
in the same specified class, then $c_u=c_v$.
\end{proposition}
\begin{proof}
For the exponential equation fix $s>0$. On $\{u<v-s\}$,
\[
 \mathrm{MA}_\theta(u)
 =\mathrm{e}^{\lambda(u-v)}\mathrm{MA}_\theta(v)
 \leq \mathrm{e}^{-\lambda s}\mathrm{MA}_\theta(v-s).
\]
Strict domination gives $u\geq v-s$. Let $s\downarrow0$ and exchange $u,v$.

For the second assertion, suppose $c_u<c_v$. For every real $A$, the global inequality
\[
 \mathrm{MA}_\theta(u)
 =\frac{c_u}{c_v}\mathrm{MA}_\theta(v+A)
\]
implies $u\geq v+A$ by strict domination. Choose a point where both potentials are finite and let $A\to+\infty$, a contradiction. Hence $c_u\geq c_v$; exchanging the two proves equality. There is no factor strictly below one when $c_u=c_v$, so this argument does not prove equality of the potentials for the constant-density equation.
\end{proof}

The corresponding uniqueness statements in the minimal ordinary type of an arbitrary smooth big background follow from \cref{cor:big-barrier-orders}. That argument uses a strict singular barrier; it does not require semipositivity or bounded mass.
